%%%%%%%%%%%%%%%%%%%%%%%%%%%%%%%%%%%%%%%%%%%%%%%%%%%%%%%%%%%%%%%%%%%%%%%%%%%%%%%
% Harvard Academic Notes - 통합 마스터 템플릿
% 모든 강의 노트에 적용되는 통일된 스타일
% 버전: 2.1 - 가독성 개선 (선택적 최적화)
% 최종 수정일: 2025-11-17
%%%%%%%%%%%%%%%%%%%%%%%%%%%%%%%%%%%%%%%%%%%%%%%%%%%%%%%%%%%%%%%%%%%%%%%%%%%%%%%

\documentclass[11pt,a4paper]{article}

%========================================================================================
% 기본 패키지
%========================================================================================

% --- 한국어 지원 ---
\usepackage{kotex}

% --- 페이지 레이아웃 ---
\usepackage[top=20mm, bottom=20mm, left=20mm, right=18mm]{geometry}
\usepackage{setspace}
\onehalfspacing                      % 1.5배 줄간격
\setlength{\parskip}{0.5em}          % 문단 간격
\setlength{\parindent}{0pt}          % 들여쓰기 없음

% --- 표 관련 ---
\usepackage{booktabs}              % 고품질 표
\usepackage{tabularx}              % 자동 너비 조절 표
\usepackage{array}                 % 표 컬럼 확장
\usepackage{longtable}             % 여러 페이지 표
\renewcommand{\arraystretch}{1.1}  % 표 행간 조절

%========================================================================================
% 헤더 및 푸터
%========================================================================================

\usepackage{fancyhdr}
\pagestyle{fancy}
\fancyhf{}
\fancyhead[L]{\small\textit{FIN-571: Financial Institutions & Monetary Policy}}
\fancyhead[R]{\small\textit{Module 06}}
\fancyfoot[C]{\thepage}
\renewcommand{\headrulewidth}{0.5pt}
\renewcommand{\footrulewidth}{0.3pt}

% 첫 페이지는 헤더 없음
\fancypagestyle{firstpage}{
    \fancyhf{}
    \fancyfoot[C]{\thepage}
    \renewcommand{\headrulewidth}{0pt}
}

%========================================================================================
% 색상 정의 (파스텔 톤 + 다크모드 호환)
%========================================================================================

\usepackage[dvipsnames]{xcolor}

% 밝은 배경용 파스텔 색상
\definecolor{lightblue}{RGB}{220, 235, 255}      % 부드러운 파랑
\definecolor{lightgreen}{RGB}{220, 255, 235}     % 부드러운 초록
\definecolor{lightyellow}{RGB}{255, 250, 220}    % 부드러운 노랑
\definecolor{lightpurple}{RGB}{240, 230, 255}    % 부드러운 보라
\definecolor{lightgray}{gray}{0.95}              % 밝은 회색
\definecolor{lightpink}{RGB}{255, 235, 245}      % 부드러운 핑크
\definecolor{boxgray}{gray}{0.95}
\definecolor{boxblue}{rgb}{0.9, 0.95, 1.0}
\definecolor{boxred}{rgb}{1.0, 0.95, 0.95}

% 진한 색상 (테두리/제목용)
\definecolor{darkblue}{RGB}{50, 80, 150}
\definecolor{darkgreen}{RGB}{40, 120, 70}
\definecolor{darkorange}{RGB}{200, 100, 30}
\definecolor{darkpurple}{RGB}{100, 60, 150}

%========================================================================================
% 박스 환경 (tcolorbox) - 6가지 타입
%========================================================================================

\usepackage[most]{tcolorbox}
\tcbuselibrary{skins, breakable}

% 1. 개요 박스 (강의 시작 부분)
\newtcolorbox{overviewbox}[1][]{
    enhanced,
    colback=lightpurple,
    colframe=darkpurple,
    fonttitle=\bfseries\large,
    title=📚 강의 개요,
    arc=3mm,
    boxrule=1pt,
    left=8pt,
    right=8pt,
    top=8pt,
    bottom=8pt,
    breakable,
    #1
}

% 2. 요약 박스
\newtcolorbox{summarybox}[1][]{
    enhanced,
    colback=lightblue,
    colframe=darkblue,
    fonttitle=\bfseries,
    title=📝 핵심 요약,
    arc=2mm,
    boxrule=0.7pt,
    left=6pt,
    right=6pt,
    top=6pt,
    bottom=6pt,
    breakable,
    #1
}

% 3. 핵심 정보 박스
\newtcolorbox{infobox}[1][]{
    enhanced,
    colback=lightgreen,
    colframe=darkgreen,
    fonttitle=\bfseries,
    title=💡 핵심 정보,
    arc=2mm,
    boxrule=0.7pt,
    left=6pt,
    right=6pt,
    top=6pt,
    bottom=6pt,
    breakable,
    #1
}

% 4. 주의사항 박스
\newtcolorbox{warningbox}[1][]{
    enhanced,
    colback=lightyellow,
    colframe=darkorange,
    fonttitle=\bfseries,
    title=⚠️ 주의사항,
    arc=2mm,
    boxrule=0.7pt,
    left=6pt,
    right=6pt,
    top=6pt,
    bottom=6pt,
    breakable,
    #1
}

% 5. 예제 박스
\newtcolorbox{examplebox}[1][]{
    enhanced,
    colback=lightgray,
    colframe=black!60,
    fonttitle=\bfseries,
    title=📖 예제,
    arc=2mm,
    boxrule=0.7pt,
    left=6pt,
    right=6pt,
    top=6pt,
    bottom=6pt,
    breakable,
    #1
}

% 6. 정의 박스
\newtcolorbox{definitionbox}[1][]{
    enhanced,
    colback=lightpink,
    colframe=purple!70!black,
    fonttitle=\bfseries,
    title=📌 정의,
    arc=2mm,
    boxrule=0.7pt,
    left=6pt,
    right=6pt,
    top=6pt,
    bottom=6pt,
    breakable,
    #1
}

% 7. 중요 박스 (importantbox - warningbox와 유사)
\newtcolorbox{importantbox}[1][]{
    enhanced,
    colback=boxred,
    colframe=red!70!black,
    fonttitle=\bfseries,
    title=⚠️ 매우 중요,
    arc=2mm,
    boxrule=0.7pt,
    left=6pt,
    right=6pt,
    top=6pt,
    bottom=6pt,
    breakable,
    #1
}

% 8. cautionbox (warningbox와 동일)
\let\cautionbox\warningbox
\let\endcautionbox\endwarningbox

\tcbset{
    summarybox/.style={enhanced, colback=lightblue, colframe=darkblue, fonttitle=\bfseries, title=핵심 요약, arc=2mm, boxrule=0.7pt, breakable},
    warningbox/.style={enhanced, colback=lightyellow, colframe=darkorange, fonttitle=\bfseries, title=주의, arc=2mm, boxrule=0.7pt, breakable},
    examplebox/.style={enhanced, colback=lightgray, colframe=black!60, fonttitle=\bfseries, title=예제, arc=2mm, boxrule=0.7pt, breakable},
    definitionbox/.style={enhanced, colback=lightpink, colframe=purple!70!black, fonttitle=\bfseries, title=정의, arc=2mm, boxrule=0.7pt, breakable},
    importantbox/.style={enhanced, colback=boxred, colframe=red!70!black, fonttitle=\bfseries, title=매우 중요, arc=2mm, boxrule=0.7pt, breakable},
    infobox/.style={enhanced, colback=lightgreen, colframe=darkgreen, fonttitle=\bfseries, title=핵심 정보, arc=2mm, boxrule=0.7pt, breakable},
    conceptbox/.style={enhanced, colback=lightgreen, colframe=darkgreen, fonttitle=\bfseries, title=개념, arc=2mm, boxrule=0.7pt, breakable},
    analogybox/.style={enhanced, colback=green!5!white, colframe=green!60!black, fonttitle=\bfseries, title=비유, arc=2mm, boxrule=0.7pt, breakable},
    overviewbox/.style={enhanced, colback=lightpurple, colframe=darkpurple, fonttitle=\bfseries\large, title=강의 개요, arc=3mm, boxrule=1pt, breakable},
    cautionbox/.style={enhanced, colback=lightyellow, colframe=darkorange, fonttitle=\bfseries, title=주의, arc=2mm, boxrule=0.7pt, breakable},
}

%========================================================================================
% 코드 블록 설정 (밝은 배경)
%========================================================================================

\usepackage{listings}

\definecolor{codegray}{rgb}{0.5,0.5,0.5}
\definecolor{codepurple}{rgb}{0.58,0,0.82}
\definecolor{backcolour}{rgb}{0.95,0.95,0.95}

\lstset{
    basicstyle=\ttfamily\small,
    backgroundcolor=\color{lightgray},
    keywordstyle=\color{darkblue}\bfseries,
    commentstyle=\color{darkgreen}\itshape,
    stringstyle=\color{purple!80!black},
    numberstyle=\tiny\color{black!60},
    numbers=left,
    numbersep=8pt,
    breaklines=true,
    breakatwhitespace=false,
    frame=single,
    frameround=tttt,
    rulecolor=\color{black!30},
    captionpos=b,
    showstringspaces=false,
    tabsize=2,
    xleftmargin=15pt,
    xrightmargin=5pt,
    escapeinside={\%*}{*)}
}

% Python 코드 스타일
\lstdefinestyle{pythonstyle}{
    language=Python,
    morekeywords={self, True, False, None},
}

% SQL 코드 스타일
\lstdefinestyle{sqlstyle}{
    language=SQL,
    morekeywords={SELECT, FROM, WHERE, JOIN, GROUP, BY, ORDER, HAVING},
}

%========================================================================================
% 목차 스타일링
%========================================================================================

\usepackage{tocloft}
\renewcommand{\cftsecleader}{\cftdotfill{\cftdotsep}}
\setlength{\cftbeforesecskip}{0.4em}
\renewcommand{\cftsecfont}{\bfseries}
\renewcommand{\cftsubsecfont}{\normalfont}

%========================================================================================
% 표 및 그림
%========================================================================================

\usepackage{graphicx}              % 이미지
\usepackage{adjustbox}             % 표/박스 크기 조절

% 표 캡션 스타일
\usepackage{caption}
\captionsetup[table]{
    labelfont=bf,
    textfont=it,
    skip=5pt
}
\captionsetup[figure]{
    labelfont=bf,
    textfont=it,
    skip=5pt
}

%========================================================================================
% 수학
%========================================================================================

\usepackage{amsmath, amssymb, amsthm}

% 정리 환경
\theoremstyle{definition}
\newtheorem{theorem}{정리}[section]
\newtheorem{lemma}[theorem]{보조정리}
\newtheorem{proposition}[theorem]{명제}
\newtheorem{corollary}[theorem]{따름정리}
\newtheorem{definition}{정의}[section]
\newtheorem{example}{예제}[section]

%========================================================================================
% 하이퍼링크
%========================================================================================

\usepackage[
    colorlinks=true,
    linkcolor=blue!80!black,
    urlcolor=blue!80!black,
    citecolor=green!60!black,
    bookmarks=true,
    bookmarksnumbered=true,
    pdfborder={0 0 0}
]{hyperref}

% PDF 메타데이터는 각 문서에서 설정
\hypersetup{
    pdftitle={FIN-571: Financial Institutions & Monetary Policy - Module 06},
    pdfauthor={강의 노트},
    pdfsubject={Academic Notes}
}

%========================================================================================
% 기타 유용한 패키지
%========================================================================================

\usepackage{enumitem}              % 리스트 커스터마이징
\setlist{nosep, leftmargin=*, itemsep=0.3em}

\usepackage{microtype}             % 타이포그래피 개선
\usepackage{footnote}              % 각주 개선
\usepackage{url}                   % URL 줄바꿈
\urlstyle{same}

%========================================================================================
% 사용자 정의 명령어
%========================================================================================

% 강조 텍스트
\newcommand{\important}[1]{\textbf{\textcolor{red!70!black}{#1}}}
\newcommand{\keyword}[1]{\textbf{#1}}
\newcommand{\term}[1]{\textit{#1}}
\newcommand{\code}[1]{\texttt{#1}}

% 용어 설명 (인라인)
\newcommand{\defterm}[2]{\textbf{#1}\footnote{#2}}

% 섹션 시작 전 페이지 분리
\newcommand{\newsection}[1]{\newpage\section{#1}}

%========================================================================================
% 문서 제목 스타일
%========================================================================================

\usepackage{titling}
\pretitle{\begin{center}\LARGE\bfseries}
\posttitle{\par\end{center}\vskip 0.5em}
\preauthor{\begin{center}\large}
\postauthor{\end{center}}
\predate{\begin{center}\large}
\postdate{\par\end{center}}

%========================================================================================
% 섹션 제목 간격
%========================================================================================

\usepackage{titlesec}
\titlespacing*{\section}{0pt}{1.5em}{0.8em}
\titlespacing*{\subsection}{0pt}{1.2em}{0.6em}
\titlespacing*{\subsubsection}{0pt}{1em}{0.5em}

%========================================================================================
% 메타 정보 박스 명령어
%========================================================================================

\newcommand{\metainfo}[4]{
\begin{tcolorbox}[
    colback=lightpurple,
    colframe=darkpurple,
    boxrule=1pt,
    arc=2mm,
    left=10pt,
    right=10pt,
    top=8pt,
    bottom=8pt
]
\begin{tabular}{@{}rl@{}}
▣ \textbf{강의명:} & #1 \\[0.3em]
▣ \textbf{주차:} & #2 \\[0.3em]
▣ \textbf{교수명:} & #3 \\[0.3em]
▣ \textbf{목적:} & \begin{minipage}[t]{0.75\textwidth}#4\end{minipage}
\end{tabular}
\end{tcolorbox}
}

%========================================================================================
% 끝
%========================================================================================


\begin{document}

\thispagestyle{firstpage}

\metainfo{중앙은행과 통화정책}{모듈 2 (파트 1/2)}{Prof. Ralf Meisenzahl}{미국 연방준비제도(Fed)의 역사, 구조, 기능 및 통화정책 도구를 심층적으로 학습하고, 연준의 이중 책무(최대 고용, 물가 안정)를 중심으로 전통적 및 비전통적 통화정책 수단들을 이해한다.}

\tableofcontents

\newpage

% 본문 시작
이 문서는 '중앙은행과 통화정책' 강의의 모듈 2 중 첫 번째 파트 내용을 담고 있습니다.
미국 연방준비제도(Federal Reserve System)의 역사, 구조, 기능, 그리고 통화정책 도구들을 심층적으로 다룹니다.
특히, 연준의 이중 책무(Dual Mandate)인 최대 고용과 물가 안정을 중심으로 전통적 및 비전통적 통화정책 수단들을 상세히 살펴봅니다.
화폐 공급의 개념과 측정, 그리고 연준의 대차대조표 변화를 통해 통화정책이 경제에 미치는 영향을 이해하는 것이 목표입니다.
이 문서는 비전공자나 처음 배우는 사람도 완벽히 이해할 수 있도록 쉬운 설명과 풍부한 예시를 제공합니다.

\vspace{1em}

\begin{summarybox}
\textbf{문서 전체 핵심 요약:}
본 문서는 미국 연방준비제도(Fed)의 역사, 구조, 주요 기능 및 통화정책 도구를 다룹니다.
연준의 이중 책무(최대 고용, 물가 안정)를 중심으로 전통적 통화정책(공개시장운영, 재할인 창구)과 비전통적 통화정책(선제 안내, 대규모 자산 매입)을 설명합니다.
화폐 공급의 개념과 측정, 그리고 연준 대차대조표의 변화를 통해 통화정책의 작동 원리를 이해하는 데 중점을 둡니다.
\end{summarybox}

\section{용어 정리}

다음은 본 문서에서 다루는 주요 용어들을 비전공자도 이해하기 쉽게 정리한 표입니다.

\begin{adjustbox}{width=\textwidth,center}
\begin{tabular}{lllc}
\toprule
\textbf{용어 (Term)} & \textbf{쉬운 설명 (Easy Explanation)} & \textbf{원어 (Original Term)} & \textbf{비고 (Notes)} \\
\midrule
중앙은행 & 한 국가의 금융 시스템을 관리하고 통화 정책을 수행하는 기관 & Central Bank & 스웨덴 릭스방크, 미국 연방준비제도 등 \\
연방준비제도 & 미국의 중앙은행 시스템 & Federal Reserve System (Fed) & 줄여서 '연준'이라고도 함 \\
이중 책무 & 연준의 두 가지 주요 목표: 최대 고용과 물가 안정 & Dual Mandate & 통화 정책의 핵심 목표 \\
최대 고용 & 일자리를 찾는 대부분의 사람들이 일하고 있는 상태 & Maximum Employment & 실업률이 낮은 상태를 의미 \\
물가 안정 & 물가 상승률(인플레이션)이 낮고 안정적인 상태 & Stable Prices & 저축과 투자를 장려 \\
통화 정책 & 중앙은행이 경제의 화폐 공급과 신용 조건을 조절하는 활동 & Monetary Policy & 금리 조절, 화폐 공급 조절 등 \\
지급결제 시스템 & 돈이 한 사람/기관에서 다른 사람/기관으로 안전하게 이동하는 방식 & Payment System & 수표, ACH, 대규모 거래 등 \\
화폐 공급 & 경제 내에 유통되는 돈의 총량 & Money Supply & 통화량이라고도 함 \\
통화 승수 & 은행의 대출 활동을 통해 화폐 공급이 얼마나 증가하는지 나타내는 비율 & Money Multiplier & 예금의 일부를 대출하여 새로운 예금을 창출 \\
본원 통화 & 중앙은행이 직접 발행하는 화폐 (현금 + 은행의 중앙은행 예치금) & Monetary Base & 가장 좁은 의미의 화폐 공급 \\
M1 & 본원 통화에 요구불 예금(수표 계좌 잔액)을 더한 화폐 공급 지표 & M1 & 유동성이 높은 자산 포함 \\
M2 & M1에 저축성 예금, 소액 정기 예금 등을 더한 화폐 공급 지표 & M2 & 가장 일반적으로 사용되는 지표 \\
M3 & M2에 대규모 정기 예금, 기관 머니마켓 펀드 등을 더한 화폐 공급 지표 & M3 & 현재 연준에서 발표하지 않음 \\
양적 화폐 이론 & 화폐 공급량과 물가 수준, 경제 성장 간의 관계를 설명하는 이론 & Quantity Theory of Money & 화폐 공급 증가가 인플레이션을 유발할 수 있음 \\
공개시장운영 & 중앙은행이 국채 등을 사고팔아 화폐 공급과 금리를 조절하는 정책 & Open Market Operations (OMO) & 전통적인 통화 정책 수단 \\
연방기금 금리 & 은행들이 서로 하룻밤 동안 돈을 빌려줄 때 적용되는 금리 & Federal Funds Rate & 연준의 주요 목표 금리 \\
재할인 창구 & 중앙은행이 은행에 단기 대출을 제공하는 제도 & Discount Window & 비상시 유동성 공급 수단 \\
선제 안내 & 중앙은행이 미래의 통화 정책 방향에 대해 시장에 미리 알려주는 것 & Forward Guidance & 장기 금리 목표에 사용 \\
대규모 자산 매입 & 중앙은행이 국채나 주택저당증권 등을 대규모로 매입하는 정책 & Large-Scale Asset Purchases (LSAP) & 장기 금리를 직접 낮추는 효과 \\
오버나이트 역환매 조건부 & 중앙은행이 금융기관에 단기 증권을 팔고 다음날 다시 사는 거래 & Overnight Reverse Repurchase Program (ON RRP) & 단기 금리 조절에 사용 \\
대차대조표 & 특정 시점의 자산, 부채, 자본 상태를 보여주는 재무 보고서 & Balance Sheet & 연준의 자산과 부채 현황 \\
자산 & 경제적 가치가 있는 것 (예: 연준이 보유한 국채) & Assets & 미래 경제적 효익을 제공 \\
부채 & 갚아야 할 의무 (예: 시중에 유통되는 화폐, 은행의 중앙은행 예치금) & Liabilities & 미래 경제적 효익을 포기 \\
\bottomrule
\end{tabular}
\end{adjustbox}

\newpage
\section{모듈 2: 연방준비제도 개요}

안녕하세요, 여러분. 어디에서 참여하시든 잘 지내시길 바랍니다.

이번 모듈에서는 미국 연방준비제도(Federal Reserve System), 즉 연준의 기능과 통화정책 도구들을 심층적으로 학습합니다. 우리는 연준의 역사부터 시작하여, 왜 중앙은행이 필요한지, 어떻게 조직되어 있는지, 그리고 어떤 목적을 수행하는지 살펴볼 것입니다.

\begin{conceptbox}
\textbf{연준의 이중 책무 (Dual Mandate)}
연준의 통화정책 목표는 크게 두 가지입니다: \textbf{최대 고용(Maximum Employment)}과 \textbf{물가 안정(Stable Prices)}. 이 두 가지는 종종 연준의 '이중 책무'라고 불리며, 이 강의의 모든 논의는 이 이중 책무와 연결됩니다.
\end{conceptbox}

우리는 연준의 제도적 특징과 통화정책 의사결정 과정이 어떻게 작동하는지 면밀히 검토할 것입니다. 명확히 하자면, 이 강의는 미국의 연방준비제도 시스템을 사례 연구로 다룹니다. 물론 전 세계에는 많은 다른 중앙은행들이 존재하며, 일부 중앙은행은 다른 책무를 가지거나 약간 다른 핵심 기능을 수행하기도 합니다. 그러나 모든 중앙은행은 통화정책을 수행하고 지급결제 시스템을 제공합니다. 이 강의에서는 주로 중앙은행이 통화정책을 어떻게 수행하는지에 초점을 맞춥니다.

\vspace{1em}

\begin{summarybox}
\textbf{통화정책의 진화: 화폐 공급 vs. 금리}
시간이 지남에 따라 통화정책을 수행하는 최선의 방법에 대한 많은 이견이 있었습니다. 일부 중앙은행은 금리보다는 화폐 공급을 명시적으로 목표로 삼기도 했습니다. 연준도 과거에는 화폐 공급을 목표로 삼았습니다.
\end{summarybox}

\textbf{왜 화폐 공급에 관심을 가져야 할까요?}
다른 상품과 마찬가지로, 돈의 가격은 경제 내의 돈에 대한 수요와 공급에 의해 결정됩니다. 아마도 특이한 점은 돈의 가격이 바로 금리라는 것입니다. 금리를 직접 목표로 삼기보다는 화폐 공급을 통제하는 것이 더 합리적일 수 있습니다. 하지만 이는 화폐 공급을 어떻게 측정할 수 있는가라는 질문을 던집니다. 따라서 우리는 화폐 공급의 개념, 그 측정 방법, 그리고 통화정책에 의해 어떻게 영향을 받을 수 있는지 소개할 것입니다.

그 다음으로, 우리는 연준의 대차대조표와 그것이 시간이 지남에 따라 어떻게 변화해왔는지 살펴볼 것입니다. 통화정책 수행이 화폐 공급을 어떻게 변화시켰는지에 대한 통찰력을 제공할 것입니다.

\vspace{1em}

\begin{conceptbox}
\textbf{전통적 통화정책 (Traditional Monetary Policy)}
화폐 공급과 금리가 정확히 어떻게 관련되어 있는지 이해하기 위해, 우리는 전통적인 통화정책에 초점을 맞출 것입니다.
\begin{itemize}
    \item \textbf{공개시장운영 (Open Market Operations):} 공개시장운영이 연방기금 시장(federal funds market)에서 화폐 공급에 직접적으로 어떻게 영향을 미치는지 논의할 것입니다. 이러한 공개시장운영이 이 시장의 균형 금리인 \textbf{연방기금 금리(Federal funds rate)}를 어떻게 결정하는지 살펴볼 것입니다. 연방기금 금리는 미국 통화정책의 주요 목표 금리입니다.
    \item \textbf{재할인 창구 (Discount Window):} 보완적인 전통적 통화정책 도구는 재할인 창구입니다. 이 정책 도구는 은행에 단기 대출 형태로 단기 유동성을 제공합니다. 이 도구는 주로 비상용으로 설계되었기 때문에 금리는 연방기금 금리보다 높게 설정됩니다. 그러나 재할인 창구 대출은 연방기금의 대안으로 남아 있습니다. 따라서 재할인 창구 대출의 금리는 연준이 설정한 연방기금 금리 목표 범위의 상한선을 강제합니다.
\end{itemize}
\end{conceptbox}

\vspace{1em}

\begin{conceptbox}
\textbf{비전통적 통화정책 (Non-traditional Monetary Policy)}
이 모듈의 마지막 부분에서는 2008년 글로벌 금융 위기 이후 등장한 새로운 통화정책 도구들을 살펴봅니다. 이들은 종종 '비전통적 통화정책 도구'라고 불립니다.
\begin{itemize}
    \item \textbf{선제 안내 (Forward Guidance):} 선제 안내는 시장 참여자들에게 미래의 금리 또는 더 일반적으로 통화정책 조치에 대한 명확한 경로를 제공하며, 장기 금리를 목표로 하는 데 사용될 수 있습니다.
    \item \textbf{대규모 자산 매입 (Large-Scale Asset Purchases):} 단기 금리가 이미 0일 때, 중앙은행은 국채나 주택저당증권 등을 대규모로 매입하여 화폐 공급을 더욱 늘릴 수 있습니다. 이 접근 방식은 중앙은행이 장기 금리를 직접 낮출 수 있도록 합니다.
    \item \textbf{오버나이트 역환매 조건부 프로그램 (Overnight Reverse Repurchase Program, ON RRP):} 이 도구는 환매 시장(repo market)에서 금리를 설정하는 데 사용됩니다. 환매 시장은 연방기금 시장과 유사한 오버나이트 대출 시장입니다. 금융 기관들은 이 시장들 사이를 전환할 수 있으며, 이들은 대체재입니다. 환매 시장의 균형을 관리함으로써 연준은 모든 단기 금리가 통화정책 목표와 일치하도록 보장합니다.
\end{itemize}
\end{conceptbox}

결론적으로, 이 모듈은 통화정책 도구에 대한 심층적인 논의를 제공합니다. 이 비디오들 중 다수는 첫 번째 모듈에서 논의된 금융 시장 및 상품에 대한 이해를 필요로 할 것입니다. 특히, 연방기금 시장, 환매 시장, 국채, 그리고 수익률 곡선에 대한 이해가 중요합니다.

\begin{summarybox}
\textbf{핵심 요약: 통화정책 도구의 두 가지 주요 관점}
모든 통화정책 도구의 기술적인 세부 사항을 모두 이해할 필요는 없습니다. 그러나 다음 두 가지 주요 고수준 포인트를 명심하시기 바랍니다.
\begin{enumerate}
    \item \textbf{다중 시장 목표:} 연방준비제도는 금리 정책을 통해 여러 밀접하게 연결된 시장을 목표로 합니다.
    \item \textbf{제로 금리 하한선 극복:} 단기 금리가 0일 때도 중앙은행은 선택지가 없는 것이 아닙니다. 선제 안내와 대규모 자산 매입을 사용하여 장기 금리를 목표로 할 수 있습니다.
\end{enumerate}
\end{summarybox}

\newpage
\subsection{레슨 2-1.1. 연방준비제도의 역사와 구조}

안녕하세요, 연방준비제도 시스템의 역사와 구조에 대한 강의에 오신 것을 환영합니다. 이 수업에서는 중앙은행의 개념을 소개할 것입니다. 중앙은행은 전 세계에 존재합니다. 우리는 중앙은행이 왜 설립되었는지, 그들의 목표와 주요 기능은 무엇인지 논의할 것입니다. 그런 다음 미국의 중앙은행인 연방준비제도 시스템에 초점을 맞출 것입니다. 우리는 그 설립을 논의하고, 연방준비제도 시스템의 구조와 의사결정 기구를 검토할 것입니다.

\vspace{1em}

\subsubsection*{세계에서 가장 오래된 중앙은행: 스웨덴 릭스방크 사례}
세계에서 가장 오래된 중앙은행인 스웨덴 릭스방크(Swedish Riksbank)부터 살펴보겠습니다. 스웨덴 스톡홀름에 기반을 둔 이 은행은 국가들이 왜 중앙은행을 설립했는지, 그리고 무엇을 달성하고자 했는지 이해하는 데 도움이 되는 흥미로운 사례 연구입니다.

\begin{examplebox}
\textbf{스웨덴 릭스방크의 설립 배경 (1661년)}
\begin{itemize}
    \item \textbf{초기 화폐:} 1661년 스웨덴에서는 귀금속으로 만든 무거운 동전이 화폐로 사용되었습니다. 이 동전들은 무게가 45파운드(약 20kg)에 달하여 매우 불편했습니다.
    \item \textbf{지폐의 등장:} 1661년, 스웨덴 국왕은 민간 은행들이 종이 은행권을 발행하도록 허용했습니다. 이것이 오늘날 우리가 아는 지폐의 시작입니다. 지폐는 은행과 시민들에게 훨씬 더 편리했습니다.
    \item \textbf{지폐의 신뢰성:} 사람들은 왜 이 은행권을 기꺼이 받아들였을까요? 은행이 예금된 동전과 은행권을 교환해주겠다고 약속했기 때문입니다. 은행권은 요구불 상환(redeemable on demand)이 가능했습니다.
    \item \textbf{은행권 남발과 뱅크런:} 은행권은 큰 성공을 거두었지만, 은행은 점점 더 많은 은행권을 발행했습니다. 이로 인해 사람들은 민간 은행이 모든 은행권을 교환해줄 만큼 충분한 동전을 금고에 보관하고 있는지 의문을 품기 시작했습니다. 이러한 의심은 은행권을 동전으로 상환하려는 사람들이 늘어나게 만들었습니다. 이것이 바로 \textbf{뱅크런(bank run)}의 한 예시입니다. 결국 은행은 모든 은행권을 한 번에 상환할 만큼 충분한 동전을 가지고 있지 않았고, 그 결과 은행은 파산했습니다.
\end{itemize}
\end{examplebox}

\vspace{1em}

이 뱅크런에도 불구하고, 지폐의 아이디어는 여전히 인기가 있었습니다. 그러나 예금된 동전에 비해 너무 많은 지폐가 유통되고 있다는 점이 문제로 인식되었습니다. 즉, 은행 시스템이 보유한 담보(collateral)에 비해 너무 많은 돈이 유통되고 있었던 것입니다.

\begin{conceptbox}
\textbf{스웨덴 릭스방크의 설립 (1668년) 및 중앙은행의 두 가지 주요 역할}
이러한 문제 해결을 위해 1668년 스웨덴 릭스방크가 설립되어 화폐 공급을 통제하게 되었습니다. 법은 이 새로운 중앙은행이 국내 주화의 가치를 올바르고 공정하게 유지하도록 요구했습니다. 또한 법은 릭스방크가 인플레이션을 유발할 정도로 많은 지폐를 발행하는 것을 금지했습니다.

스웨덴 릭스방크의 설립 이야기는 현대 중앙은행이 해결하도록 설계된 두 가지 중요한 문제를 보여줍니다.
\begin{enumerate}
    \item \textbf{물가 안정 보장 (Ensuring Price Stability):} 중앙은행은 일반적으로 화폐 공급에 대한 통제권을 부여받습니다. 이 통제권은 중앙은행이 너무 많은 지폐가 인쇄되는 것을 막을 수 있도록 합니다. 경제에 돈이 너무 많으면 물가 인플레이션을 유발하는 경향이 있습니다.
    \item \textbf{최후의 대부자 역할 (Lender of Last Resort):} 은행은 두려움에 찬 고객들의 뱅크런에 취약합니다. 뱅크런에 대처하기 위해 중앙은행은 '최후의 대부자' 역할을 합니다.
\end{enumerate}
\end{conceptbox}

\textbf{최후의 대부자 역할이란 무엇일까요?}
은행은 단기로 차입하여 장기로 대출합니다. 스웨덴 사례에서 은행은 은행권과 교환하여 금속 동전을 차입했습니다. 그리고 그 동전의 일부를 대출하는 데 사용했습니다. 어느 시점에서든 금고에 있는 동전은 발행된 은행권보다 적었습니다. 모든 사람이 은행권을 동전으로 바꾸고 싶어 할 때, 은행은 동전이 부족하여 파산할 수 있었습니다. 현대 은행업에서는 은행이 예금자로부터 현금을 차입하여 장기로 대출합니다. 이들도 수익성 있는 대출을 좋은 차입자에게 하고 있더라도 때때로 현금 부족에 직면할 수 있습니다.

\begin{examplebox}
\textbf{최후의 대부자 기능의 실제 작동 방식}
중앙은행의 최후의 대부자 기능은 건강한 은행이 현금 부족으로 인해 파산하는 것을 막기 위해 설계되었습니다. 이것이 실제로 어떻게 작동할까요?
\begin{itemize}
    \item \textbf{단기 대출 제공:} 중앙은행은 담보(collateral)를 받고 은행에 단기 대출 형태로 현금을 제공합니다.
    \item \textbf{비상시 사용:} 이는 은행이 스트레스를 받는 비상시에 발생할 수 있습니다.
    \item \textbf{다양한 담보 인정:} 중앙은행은 일반적으로 다양한 유형의 담보를 받습니다. 중요한 것은, 이는 유동성이 낮은 장기 대출, 즉 판매하기 어렵거나 큰 할인 없이는 판매할 수 없는 자산일 수 있다는 점입니다. 예를 들어, 오늘날 연방준비제도 시스템은 부동산 및 소비자 대출을 은행에 대한 현금 대출의 담보로 인정합니다.
\end{itemize}
\end{examplebox}

\subsubsection*{미국 연방준비제도의 설립}
이제 연방준비제도 시스템이 왜 설립되었는지 살펴보겠습니다. 연방준비제도 시스템은 미국 최초의 중앙은행이 아니며, 실제로는 세 번째입니다. 1791년부터 1811년까지 존재했던 제1차 미국 은행(First Bank of America)과 1816년부터 1836년까지 존재했던 제2차 미국 은행(Second Bank of America)이 있었습니다. 그러나 이 두 초기 중앙은행은 오늘날 우리가 중앙은행이라고 생각하는 것과는 매우 달랐습니다. 예를 들어, 이들은 화폐 공급에 대한 통제권이 없었습니다.

미국은 1837년부터 1913년까지 중앙은행이 없었습니다. 스웨덴 사례와 마찬가지로 돈은 민간에서 발행되었습니다. 이 기간 동안 1907년의 공황(panic)을 포함하여 여러 차례의 은행 위기가 발생했습니다. 1907년 공황 동안 주식 시장은 가치의 50\%를 잃었습니다. 전국적으로 은행들은 고객들이 예금을 인출하기 위해 줄을 서는 것을 목격했습니다.

\begin{conceptbox}
\textbf{1907년 공황과 연방준비제도법 (Federal Reserve Act)}
1907년 공황은 최후의 대부자의 필요성을 명확히 보여주었습니다. 그러나 연방준비제도 시스템이 공식적으로 설립되기까지는 시간이 걸렸고 많은 정치적 논쟁이 있었습니다. 1913년 12월 24일, 우드로 윌슨 대통령은 연방준비제도법에 서명했습니다.

\textbf{연방준비제도법이 연준에 허용한 것:}
\begin{enumerate}
    \item \textbf{화폐 공급 통제권 부여:} 연준은 화폐 공급에 대한 통제권을 갖게 되었습니다.
    \item \textbf{은행에 대한 최후의 대부자 역할 수행:} 연준은 은행에 대한 최후의 대부자 역할을 할 수 있게 되었습니다. 이 최후의 대부자 기능은 2008년 금융 위기와 COVID-19 팬데믹을 포함하여 시장 전반의 스트레스 시기에 개입하는 데 사용되었습니다.
\end{enumerate}
\end{conceptbox}

\subsubsection*{연방준비제도의 구조}
연방준비제도법은 새로운 중앙은행의 세 가지 부분을 설립했습니다.

\begin{enumerate}
    \item \textbf{이사회 (Board of Governors):} 워싱턴 D.C.에 위치한 독립적인 연방 기관입니다. 이사회는 7명의 이사로 구성됩니다. 이사회는 전체 연방준비제도 시스템의 정책을 안내합니다.
    \item \textbf{12개 지역 연방준비은행 (Regional Federal Reserve Banks):}
        \begin{itemize}
            \item \textbf{위치:} 보스턴, 뉴욕, 필라델피아, 클리블랜드, 리치먼드, 애틀랜타, 시카고, 세인트루이스, 미니애폴리스, 캔자스시티, 댈러스, 샌프란시스코에 위치합니다.
            \item \textbf{설립 목적:} 지역 연방준비은행들은 중요한 금융 중심지에 배치되어 연준의 권한을 미국 전역에 분산시킵니다.
            \item \textbf{주요 기능:} 이들은 연방준비제도의 정책과 기능을 실행합니다. 예를 들어, 금융 기관을 검사하고 감독하며, 해당 지역 은행에 대한 최후의 대부자 역할을 합니다. 또한 해당 지역 은행에 미국 지급결제 시스템 서비스를 제공합니다.
        \end{itemize}
    \item \textbf{연방공개시장위원회 (Federal Open Market Committee, FOMC):}
        \begin{itemize}
            \item \textbf{구성:} FOMC는 12명의 위원으로 구성됩니다. 항상 이사회 7명의 이사와 뉴욕 연방준비은행 총재가 포함됩니다. 나머지 4명의 위원은 나머지 연방준비은행 총재들의 순환 그룹에서 나옵니다.
            \item \textbf{주요 임무:} FOMC의 주요 임무는 금리 목표를 설정하는 것입니다. 이는 일반적으로 통화정책이라고 불립니다.
        \end{itemize}
\end{enumerate}

\begin{summarybox}
\textbf{레슨 2-1.1 핵심 요약}
\begin{itemize}
    \item 중앙은행이 존재하기 전에는 지폐가 종종 민간에서 발행되었습니다. 이 시스템에서는 화폐 공급이 예측 불가능했고 뱅크런이 자주 발생했습니다.
    \item 화폐 공급을 통제하고 은행 시스템의 안정성을 보장하기 위해 중앙은행이 설립되었습니다.
    \item 중앙은행은 위기 시 은행에 현금을 대출해주는 최후의 대부자 역할도 합니다.
    \item 미국에서는 연방준비제도 시스템이 이사회, 12개 연방준비은행, 연방공개시장위원회라는 세 가지 기관을 통해 이러한 기능을 수행합니다.
\end{itemize}
\end{summarybox}

\newpage
\subsection{레슨 2-1.2. 연방준비제도의 목적과 기능}

안녕하세요, 연방준비제도 시스템의 다섯 가지 핵심 기능에 대한 강의에 오신 것을 환영합니다. 이 수업에서는 연방준비제도 시스템의 책무에 대해 논의할 것입니다. 법에 따라 연준은 미국 경제의 효과적인 운영을 촉진하기 위해 존재합니다. 이것이 실제로 무엇을 의미하는지, 그리고 연준의 다섯 가지 핵심 기능이 무엇인지 살펴볼 것입니다.

\begin{conceptbox}
\textbf{연준의 다섯 가지 핵심 기능}
\begin{enumerate}
    \item 통화정책 수행 (Conduct monetary policy)
    \item 금융 시스템 안정성 증진 (Promote financial system stability)
    \item 금융 기관 규제 및 감독 (Regulate and supervise financial institutions)
    \item 효율적인 지급결제 시스템 제공 (Provide an efficient payment system)
    \item 소비자 보호 및 지역 사회 개발 증진 (Promote consumer protection and community development)
\end{enumerate}
우리는 이 기능들을 하나씩 논의하고, 이들이 연준의 목표와 어떻게 연결되는지 살펴볼 것입니다.
\end{conceptbox}

\subsubsection*{1. 통화정책 수행 (Conduct Monetary Policy)}
1913년 연방준비제도법은 연방준비제도 시스템을 창설했습니다. 이 법에 따라 연준은 생산적이고 안정적인 경제를 육성하기 위해 통화정책을 수행해야 합니다. 통화정책은 시간이 지남에 따라 진화했으며, 오늘날에는 단기 금리 통제를 포함합니다.

\begin{conceptbox}
\textbf{연준의 법적 통화정책 목표 (이중 책무 + 장기 금리)}
법에 따라 연준은 통화정책을 사용하여 다음 목표를 효과적으로 촉진해야 합니다.
\begin{itemize}
    \item \textbf{최대 고용 (Maximum Employment):} 일반적으로 일자리를 찾는 대부분의 사람들이 유급 고용되어 있는 상황을 의미합니다.
    \item \textbf{물가 안정 (Stable Prices):} 물가가 안정적일 때 저축과 투자가 장려됩니다. 반대로 인플레이션이 높으면 자산 수익률이 낮아질 수 있으며, 이는 가계의 저축 유인과 기업의 투자 유인을 감소시킵니다.
    \item \textbf{온건한 장기 금리 (Moderate Long-Term Interest Rates):} 온건한 장기 금리는 기업의 투자 유인을 증가시킵니다.
\end{itemize}
\textbf{목표 달성:} 최대 고용, 물가 안정, 합리적인 장기 금리 하에서 안정적인 성장과 번성하는 경제가 있을 것이라는 믿음이 있습니다. 이것이 연준이 달성하고자 하는 바입니다.
\end{conceptbox}

\subsubsection*{2. 금융 시스템 안정성 증진 (Promote Financial System Stability)}
금융 안정성이란 무엇일까요? 일반적으로 우리는 금융 시스템이 가계와 기업에 합리적인 비용으로 중단 없이 서비스를 제공할 때 안정적이라고 간주합니다. 금융 시스템은 금융 시장과 은행과 같은 금융 기관으로 구성됩니다. 안정적인 시스템은 저축자로부터 차입자에게 자금을 전달하며, 이는 경제 전반에 걸쳐 자본의 효율적인 배분을 가능하게 합니다.

\begin{conceptbox}
\textbf{금융 시스템 모니터링 접근 방식}
연준은 이 책임을 다하기 위해 금융 시스템을 모니터링합니다. 실제로 두 가지 접근 방식을 사용하여 모니터링합니다.
\begin{enumerate}
    \item \textbf{미시 건전성 접근 방식 (Micro-prudential Approach):} 특정 기관(은행 또는 보험 회사 등)의 위험에 초점을 맞춥니다. 이 미시 건전성 접근 방식은 은행 감독과 관련이 있으며, 이에 대해서는 잠시 후에 논의할 것입니다.
    \item \textbf{거시 건전성 접근 방식 (Macro-prudential Approach):} 대신 전체 금융 시스템에 영향을 미칠 수 있는 위험을 찾습니다. 연방준비제도는 여러 금융 시장 스트레스 측정 지표를 모니터링합니다.
\end{enumerate}
\end{conceptbox}

\begin{examplebox}
\textbf{거시 건전성 모니터링 예시}
\begin{itemize}
    \item \textbf{레버리지 (Leverage):} 금융 시스템 내의 차입 또는 레버리지의 양은 손실을 증폭시키고 채무 불이행을 유발할 수 있습니다.
    \item \textbf{단기 차입 (Short-Term Borrowing):} 금융 기업의 단기 차입량은 신용 흐름을 방해하는 뱅크런을 유발할 수 있습니다.
\end{itemize}
\end{examplebox}

\subsubsection*{3. 금융 기관 규제 및 감독 (Regulate and Supervise Financial Institutions)}
연방준비제도의 세 번째 핵심 기능은 금융 기관을 규제하고 감독하는 것입니다. 규제와 감독은 별개이지만 상호 보완적인 활동입니다.

\begin{conceptbox}
\textbf{규제 (Regulation)와 감독 (Supervision)의 차이}
\begin{itemize}
    \item \textbf{규제 (Regulation):} 규칙을 작성하는 것과 관련이 있습니다. 이러한 규칙은 종종 금융 기관에 제한을 가합니다. 이 과정은 의회의 입법으로 시작되며, 이는 연방준비제도에 규칙을 초안할 권한을 부여합니다. 대중의 의견을 수렴하여 규칙은 최종 확정되고 공표되며, 금융 기관은 이 규칙을 준수해야 합니다.
    \item \textbf{감독 (Supervision):} 반면에 규제 준수와 관련이 있습니다. 이를 위해 연방준비제도는 은행 검사관을 고용하고 훈련시킵니다. 검사관은 규제 대상 금융 기관을 방문하여 규칙 준수 여부를 확인하며, 이는 은행 운영에 대한 심층적인 검사를 포함합니다. 규칙이 위반될 경우, 감독관은 벌금 및 배당금 지급 제한과 같은 다른 제재를 부과할 수 있습니다.
\end{itemize}
\end{conceptbox}

\begin{examplebox}
\textbf{규제 예시: 진실한 대출법 (Truth in Lending Act)}
진실한 대출법은 좋은 예시입니다. 이 규칙에 따라 연방준비제도는 대출 기관이 대출 조건과 비용을 차입자에게 명확하게 공개하도록 요구합니다. 여기에는 신용 카드 이자율이 포함됩니다.

\textbf{감독의 중요성:}
감독은 금융 안정성과 관련이 있습니다. 이는 개별 금융 기관의 안전성과 건전성을 보장하는 미시 건전성 접근 방식입니다.
\end{examplebox}

\subsubsection*{4. 효율적인 지급결제 시스템 제공 (Provide an Efficient Payment System)}
연방준비제도의 네 번째 핵심 기능은 효율적인 지급결제 시스템을 제공하는 것입니다. 현재 연방준비제도는 다양한 지급결제를 지원합니다.

\begin{examplebox}
\textbf{연준이 지원하는 지급결제 시스템}
\begin{itemize}
    \item \textbf{통화 (Currency):} 연방준비제도는 전통적으로 많은 가계 거래에 사용되었던 지폐 또는 현금을 발행합니다.
    \item \textbf{수표 청산 (Check Clearing):} 연방준비제도는 수표 청산을 용이하게 합니다.
    \item \textbf{자동 어음 교환소 (Automated Clearinghouse, ACH):} ACH는 직원에게 지급되는 임금이나 공과금 납부와 같은 소액 은행 이체를 결제합니다.
    \item \textbf{대규모 거래 (Large Value Transactions):} 기업과 은행 간의 대규모 거래 또는 도매 지급결제도 지원합니다.
\end{itemize}
\end{examplebox}

\subsubsection*{5. 소비자 보호 및 지역 사회 개발 증진 (Promote Consumer Protection and Community Development)}
연방준비제도의 다섯 번째 핵심 기능은 소비자 보호 및 지역 사회 개발을 증진하는 것입니다. 저렴한 금융 상품에 대한 접근성을 높이는 것은 가계의 금융 안정성을 제공합니다. 이를 위해 연방준비제도는 다양한 소비자 보호, 공정 대출, 공정 주택 및 지역 사회 재투자 법률을 시행합니다.

\begin{examplebox}
\textbf{소비자 보호 및 지역 사회 개발 관련 법률}
\begin{itemize}
    \item \textbf{진실한 대출법 (Truth in Lending Act):} 앞서 언급했듯이, 이는 투명성과 공정성에 관한 것입니다.
    \item \textbf{지역 사회 재투자법 (Community Reinvestment Act):} 또 다른 중요한 예시입니다. 이 규칙에 따라 연방준비제도는 지역 사회로 유입되는 신용의 양에 대한 데이터를 수집하고, 금융 기관이 서비스하는 전체 지역 사회에 신용 접근성을 제공하도록 장려합니다. 여기에는 저소득 및 중간 소득 지역에 대한 대출이 포함됩니다.
\end{itemize}
\end{examplebox}

\begin{summarybox}
\textbf{레슨 2-1.2 핵심 요약}
이 세션에서는 연방준비제도 시스템의 목적과 기능에 대한 큰 그림을 검토했습니다. 우리는 연준의 다섯 가지 핵심 기능에 대해 논의했습니다.
\begin{itemize}
    \item 통화정책 수행
    \item 금융 시스템 안정성 증진
    \item 금융 기관 규제 및 감독
    \item 효율적인 지급결제 시스템 제공
    \item 소비자 보호 및 지역 사회 개발 증진
\end{itemize}
이러한 기능들은 총체적으로 미국 경제를 지원하고 더 일반적으로 공익에 기여합니다.
\end{summarybox}

\newpage
\section{레슨 2-2: 전통적 통화정책}

\subsection{레슨 2-2.1. 화폐 공급과 통화정책}

안녕하세요, 화폐 공급에 대한 강의에 오신 것을 환영합니다. 이 수업에서는 화폐 공급이 무엇이며 어떻게 측정되는지 논의할 것입니다. 우리는 화폐 공급에 대한 두 가지 아이디어를 소개할 것입니다. 첫째, 통화 승수(money multiplier), 둘째, 양적 화폐 이론(quantity theory of money)입니다. 이러한 아이디어들은 화폐 공급을 인플레이션을 포함한 중요한 경제 결과와 연결합니다. 중앙은행이 인플레이션을 목표로 한다면, 화폐 공급을 목표로 하는 것이 합리적일 수 있습니다. 우리는 중앙은행이 화폐 공급을 어떻게 통제하는지 논의할 것입니다.

\begin{conceptbox}
\textbf{화폐 공급의 정의}
돈은 우리가 물건을 사고팔 때 사용하는 것입니다. 넓게 말하면, 구매자와 판매자 사이의 신뢰할 수 있는 지불 수단입니다. \textbf{화폐 공급(Money Supply)}은 유통되는 돈의 총량입니다. 화폐 공급에는 달러 지폐와 동전이 포함되지만, 다른 항목들도 포함될 수 있습니다.
\end{conceptbox}

\subsubsection*{화폐 공급의 측정}
화폐 공급을 어떻게 측정하는지 더 정확히 알아봅시다.

\begin{enumerate}
    \item \textbf{본원 통화 (Monetary Base):} 가장 좁은 정의는 본원 통화입니다. 본원 통화는 유통되는 통화량(currency in circulation)을 포함합니다. 또한 은행의 준비금 잔액(bank reserve balances)도 포함합니다. 이는 은행이 연방준비제도 계좌에 예치해 둔 예금입니다.
    \item \textbf{M1:} 화폐 공급의 다소 넓은 정의는 M1입니다. M1은 본원 통화에서 시작하여 거래 예금(transaction deposits)을 추가합니다. 여기에는 은행의 당좌 예금(checking account) 잔액이 포함됩니다. M1의 아이디어는 현금과 은행이 대출할 수 있는 자금(loanable funds)의 양을 포착하는 것입니다.
\end{enumerate}

\textbf{왜 대출 가능 자금이 그렇게 중요하며, 화폐 공급과 어떻게 관련될까요?}
다음 예시를 고려해 봅시다.

\begin{examplebox}
\textbf{통화 승수 (Money Multiplier)의 작동 원리}
\begin{itemize}
    \item 당신이 은행에 100달러를 예금했다고 가정해 봅시다.
    \item 은행은 10\%인 10달러를 연방준비제도에 준비금으로 보관합니다. (이는 의무적이거나 단순히 은행의 신중함 때문일 수 있습니다.)
    \item 은행은 나머지 90달러를 대출해 줍니다.
    \item 이 90달러는 차입자의 당좌 예금에 새로운 예금으로 나타날 것입니다.
    \item 다시 은행은 10\%인 9달러를 연방준비제도에 준비금으로 보관하고, 81달러를 대출해 줍니다.
    \item 이제 81달러는 새로운 예금으로 나타날 것입니다.
\end{itemize}
이 과정은 계속될 수 있으며, 종종 \textbf{통화 승수}라고 불립니다. 자금이 대출될 때마다 은행 예금, 그리고 따라서 화폐 공급은 증가할 것입니다. 사실, 10\%의 준비금(reserve)이 화폐 공급을 제한하는 요인입니다.
\end{examplebox}

\vspace{1em}

\begin{enumerate}
    \setcounter{enumi}{2} % Continue numbering from 2
    \item \textbf{M2:} 화폐 공급의 가장 일반적인 정의는 M2입니다. 이는 M1에 비교적 단기간에 현금으로 전환될 수 있는 다른 항목들을 더한 것입니다. 여기에는 저축성 예금(savings deposits), 10만 달러 미만의 소액 정기 예금(small time deposits), 소매 머니마켓 뮤추얼 펀드(retail money market mutual fund shares)가 포함됩니다.
    \item \textbf{M3:} 가장 넓은 정의는 M3입니다. 이는 M2에 대규모 정기 예금(large time deposits), 기관 머니마켓 펀드(institutional money market funds), 단기 환매 조건부 계약(short-term repurchase agreements), 그리고 더 큰 유동성 자산(larger liquid assets)을 더한 것입니다.
\end{enumerate}

\subsubsection*{화폐 공급과 거시 경제 결과: 양적 화폐 이론}
역사적으로 화폐 공급 측정치는 주요 거시 경제 결과와 밀접하게 관련되어 있었습니다. 과거의 높은 인플레이션 시기는 화폐 공급이 빠르게 증가했던 시기와 일치합니다. 왜 그럴까요? 다른 상품과 마찬가지로, 돈의 가치는 유통되는 돈이 과도할 경우 하락합니다.

\begin{conceptbox}
\textbf{양적 화폐 이론 (Quantity Theory of Money)}
화폐 공급, 인플레이션, 경제 성장 간의 관계를 \textbf{양적 화폐 이론}이라고 합니다. 이는 경제학에서 가장 오래된 아이디어 중 하나이며, 시카고 대학의 밀턴 프리드먼(Milton Friedman)에 의해 대중화되었습니다.

\textbf{왜 이 관계가 중요할까요?}
중앙은행은 물가 안정과 최대 고용을 목표로 합니다. 이것이 그들의 책무입니다. 역사적으로 화폐 공급 측정치는 인플레이션 및 경제 성장과 밀접하게 관련되어 있었습니다. 이러한 이유로 중앙은행은 통화정책을 결정할 때 화폐 공급 측정치를 사용해 왔습니다.
\end{conceptbox}

\textbf{데이터로 살펴본 화폐 공급과 인플레이션/GDP}
미국의 1960년대 이후 데이터를 살펴보겠습니다. 여기에 M3 화폐 공급이 있습니다. 이제 인플레이션을 측정하기 위해 소비자 물가 지수(consumer price index)를 추가합니다. 1980년대 경기 침체 이후 약간의 차이가 있지만, M3 화폐 공급이 인플레이션을 상당히 잘 추적하는 것을 볼 수 있습니다. 마찬가지로, M3는 2000년까지 명목 국내총생산(nominal cross domestic product)을 잘 추적합니다.

\subsubsection*{중앙은행의 화폐 공급 통제 방법}
원칙적으로 연방준비제도는 화폐 공급을 통제하여 그 책무를 달성할 수 있습니다. 실제로 그들은 몇 가지 방법으로 은행을 목표로 삼아 이를 수행할 수 있습니다.

\begin{enumerate}
    \item \textbf{지급준비율 변경 (Change the Reserve Requirement):} 연방준비제도는 은행의 지급준비율을 변경할 수 있습니다. 통화 승수를 상기해 봅시다. 지급준비율이 높을수록 은행은 더 많은 준비금을 보유해야 하고, 더 적은 대출을 할 수 있습니다. 다시 말해, 연방준비제도는 통화 승수를 통제할 수 있습니다.
    \item \textbf{공개시장운영 수행 (Conduct Open Market Operations):} 연방준비제도는 공개시장운영을 수행할 수 있습니다. 공개시장운영에서 연준은 은행과 증권(예: 국채)을 거래합니다. 연방준비제도가 은행에 증권을 판매하면, 은행은 대출 가능 자금(loanable funds)을 사용하여 연준에 지불해야 합니다. 결과적으로 은행은 대출할 자금이 줄어들고, 이는 통화 승수와 화폐 공급을 감소시킬 것입니다.
\end{enumerate}

\textbf{하지만 문제가 있습니다.}
불행히도 시간이 지남에 따라 화폐 공급, 인플레이션, GDP 간의 관계는 데이터에서 덜 신뢰할 수 있게 되었습니다.
\begin{itemize}
    \item \textbf{측정의 어려움:} 첫째, 이제 새로운 지급결제 기술과 가계가 현금을 보관하는 방법이 생겨 화폐 공급을 측정하기가 더 어려워졌습니다. 연방준비제도는 M3 발표를 중단했습니다.
    \item \textbf{관계의 불안정성:} 둘째, M3와 인플레이션 또는 GDP 간의 관계는 2000년 이후 안정적이지 않았습니다. 2008년 금융 위기 이후 M3와 GDP 간의 관계는 무너졌습니다. 여기에서 명확한 차이를 볼 수 있습니다. COVID-19 팬데믹 동안에도 다시 발생했습니다.
\end{itemize}
대부분의 중앙은행은 이제 통화정책 결정을 내릴 때 화폐 공급 측정치를 덜 중요하게 여깁니다.

\begin{summarybox}
\textbf{레슨 2-2.1 핵심 요약}
이 세션에서는 화폐 공급의 개념과 통화정책과의 연결성에 대해 논의했습니다.
\begin{itemize}
    \item \textbf{화폐 공급 측정:} 우리는 본원 통화(Monetary Base), M1, M2, M3 등 여러 가지 방법으로 화폐 공급을 측정할 수 있음을 논의했습니다.
    \item \textbf{통화 승수:} 통화 승수는 은행 준비금과 화폐 공급을 연결합니다.
    \item \textbf{양적 화폐 이론:} 우리는 화폐 공급, 인플레이션, GDP 간의 관계를 살펴보았습니다.
    \item \textbf{통화정책의 영향:} 통화정책이 화폐 공급을 증가시키거나 감소시킬 수 있음을 보았습니다.
\end{itemize}
\end{summarybox}

\newpage
\subsection{레슨 2-2.2. 연방준비제도의 대차대조표}

안녕하세요, 연방준비제도 대차대조표에 대한 강의에 오신 것을 환영합니다. 연방준비제도 시스템은 미국의 중앙은행입니다. 다른 금융 기관과 마찬가지로 연준도 대차대조표를 가지고 있습니다. 모든 대차대조표와 마찬가지로 자산과 부채가 있을 것입니다. 다른 점은 연준이 대차대조표를 사용하여 통화정책 목표를 달성한다는 것입니다. 예를 들어, 연준이 자산을 매입하고 대차대조표를 확장하기로 결정하는 것은 통화정책 조치입니다.

우리는 연방준비제도 대차대조표를 면밀히 살펴볼 것입니다. 우리는 그 규모를 살펴볼 것입니다. 우리는 그 구성 요소를 검토할 것입니다. 주요 자산은 무엇이며, 이러한 자산은 어떻게 조달될까요? 그런 다음 최근 몇 년 동안 연준 대차대조표의 구성이 어떻게 변했는지 살펴볼 것입니다. 위기에 대응한 극적인 통화정책 조치가 이러한 극적인 변화를 어떻게 이끌었는지 볼 것입니다. 이것은 이 수업의 중요한 주제가 될 것입니다.

\textbf{연준 대차대조표의 규모 변화}
매주 연방준비제도는 상세한 대차대조표를 발표합니다. 이 정보는 온라인에서 제공되며 시장 참여자들은 면밀히 주시합니다.

시간에 따른 연방준비제도 대차대조표의 규모를 살펴보겠습니다. 규모는 연방준비제도의 총 자산을 의미합니다.

\begin{examplebox}
\textbf{연준 대차대조표 규모의 폭발적인 성장}
\begin{itemize}
    \item \textbf{2000년대 초:} 연방준비제도는 약 8천억 달러 상당의 자산을 보유했습니다.
    \item \textbf{2008년 금융 위기:} 대차대조표는 두 배 이상 증가했습니다. 이때 연방준비제도는 몇 가지 주요 통화정책 개입을 수행했습니다. 이러한 개입은 극적인 영향을 미쳤습니다. 2009년까지 연방준비제도는 2조 달러가 조금 넘는 자산을 보유했습니다.
    \item \textbf{2009년-2014년:} 대차대조표는 다시 두 배 이상 증가하여 4조 달러를 넘어섰습니다.
    \item \textbf{COVID-19 팬데믹:} 대차대조표 규모는 다시 두 배 증가하여 2021년 여름까지 8조 달러에 도달했습니다.
\end{itemize}
\textbf{핵심 요약:} 2008년 이후 연방준비제도 대차대조표는 폭발적으로 성장했습니다. 이러한 성장은 통화정책 조치를 반영합니다.
\end{examplebox}

\subsubsection*{연준 대차대조표의 구성 요소}

\textbf{1. 2008년 이전의 대차대조표 (정상 시기)}
가장 중요한 자산의 내역부터 살펴보겠습니다. 2008년 이전의 대차대조표는 정상 시기의 연준 운영에 대해 알려줄 것입니다.

\begin{examplebox}
\textbf{2000년 1월 12일 연준 대차대조표 (자산 측면)}
\begin{itemize}
    \item \textbf{총 자산:} 약 6,300억 달러.
    \item \textbf{정부 증권 (Government securities):} 4,850억 달러 상당으로, 전체 자산의 4분의 3을 차지하는 가장 큰 자산 포지션입니다.
    \item \textbf{환매 거래 (Report transactions):} 두 번째로 큰 포지션입니다. 이는 연방준비제도가 단기적으로 현금 대출을 하는 경우입니다. 이는 일반적으로 공개시장운영의 일부입니다.
    \textbf{통화 (Currency), 금 (Gold):} 연방준비제도는 동전과 지폐를 포함한 통화와 금도 보유합니다.
\end{itemize}
\end{examplebox}

\begin{examplebox}
\textbf{2000년 1월 12일 연준 대차대조표 (부채 측면)}
\begin{itemize}
    \item \textbf{유통 통화 (Currency in circulation):} 5,900억 달러 상당으로, 가장 큰 포지션입니다. 이는 여러분의 주머니에 있는 달러 지폐와 동전입니다.
    \item \textbf{은행의 준비금 잔액 (Reserve balances of banks):} 연방준비제도에 예치된 은행의 준비금으로, 의무 준비금과 초과 준비금 모두를 포함합니다.
    \item \textbf{본원 통화 (Monetary Base):} 유통 통화와 준비금을 합친 것을 본원 통화라고 합니다. 본원 통화는 연준의 부채입니다. 이는 연준에 의해 발행되며 연방준비제도 시스템의 자산에 의해 뒷받침됩니다.
    \item \textbf{재무부 일반 계정 (Treasury General Account):} 미국 정부의 당좌 예금 계좌입니다. 모든 미국 정부 지불은 이 계좌에서 이루어집니다. 세금 수입도 이 계좌에 예치됩니다.
\end{itemize}
2000년 1월 12일의 대차대조표는 20년 후와 비교하면 작았습니다. 실제로 2000년대 초에는 COVID-19 팬데믹 기간의 대차대조표 규모의 10분의 1도 되지 않았습니다. 이러한 대차대조표 확장은 어떻게 일어났을까요?
\end{examplebox}

\textbf{2. 2008년 이후의 대차대조표 (확장 시기)}
먼저 자산 측면을 살펴보겠습니다. 다음은 2021년 1월 7일 데이터입니다. 2021년 데이터에서 주목할 점이 몇 가지 있습니다.

\begin{examplebox}
\textbf{2021년 1월 7일 연준 대차대조표 (자산 측면)}
\begin{itemize}
    \item \textbf{총 자산:} 7.3조 달러. (이전 6,300억 달러에 비해 훨씬 커졌습니다.)
    \item \textbf{항목 변화:} 주택저당증권(mortgage-backed securities)과 위기 시설(crisis facilities)이 나타납니다.
    \item \textbf{미국 국채 (US treasury securities):} 4.6조 달러 이상으로 여전히 가장 큰 항목입니다.
    \item \textbf{주택저당증권 (Mortgage-backed securities, MBS):} 정부 보증 기업(government sponsored enterprises)의 주택저당증권으로 뒷받침되는 채권으로, 2조 달러 이상을 보유했습니다. 이 채권들은 미국 정부에 의해 보증됩니다.
    \item \textbf{자본 이득 (Capital gains):} 아직 판매되지 않은 증권의 이익입니다. 이는 미실현 자본 이득(unrealized capital gains)이라고도 불립니다.
    \item \textbf{위기 시설 (Crisis facilities):} 다양한 대출 프로그램입니다. 주로 은행에 대한 것이며, 시장 전반의 스트레스 시기에 자금이 제공되었습니다. 은행은 이 대출금을 상환해야 합니다.
\end{itemize}
\textbf{자산 성장의 주요 동인:} 분명히 연방준비제도가 보유한 정부 및 정부 보증 증권입니다.
\end{examplebox}

\textbf{연준은 이 대규모 대차대조표 확장을 어떻게 조달했을까요?}
일반적인 기업이나 가계는 확장을 위해 차입해야 할 수도 있습니다. 그러나 연준은 화폐 공급을 통제하므로 다른 방식으로 이를 수행할 수 있었습니다.

\textbf{연방준비제도는 더 많은 지폐를 인쇄했을까요?}
같은 기간 동안 유통되는 총 통화량을 살펴보겠습니다. 보시다시피, 유통 통화량은 많이 증가했습니다. 20년 동안 거의 4배 증가했습니다. 그러나 이 증가가 전체 이야기는 아닙니다.

\begin{examplebox}
\textbf{2021년 1월 7일 연준 대차대조표 (부채 측면)}
\begin{itemize}
    \item \textbf{지불 방식:} 연방준비제도는 미국 재무부로부터 직접 정부 증권을 매입하는 것이 아니라 은행으로부터 매입합니다. 연준은 이 은행들에게 어떻게 지불했을까요? 지폐로 지불하지 않았습니다. 대신 디지털 화폐로 지불했습니다. 이는 연준에 있는 은행의 준비금 계좌에 대한 신용(credit) 형태로 이루어집니다.
    \item \textbf{은행 준비금 증가:} 연준은 시스템 내 은행 준비금의 양을 증가시켰습니다. 보시다시피, 연방준비제도에 있는 은행 준비금은 2021년까지 부채의 약 40\%로 확장되었습니다.
    \item \textbf{재무부 일반 계정:} 여기에 빠진 것은 무엇일까요? 나머지 부채의 대부분은 미국 정부 일반 계정(US government general account)에 속하며, 최고점에서는 약 20\%를 차지했습니다. 왜 그럴까요? 2008년 이후, 특히 2020년에 미국 정부의 역할이 커졌습니다. 이로 인해 연방준비제도에 더 많은 돈이 예치되었습니다. 정부는 이러한 예치금을 다양한 재정 프로그램에 지불하는 데 사용했습니다.
\end{itemize}
\end{examplebox}

\begin{summarybox}
\textbf{레슨 2-2.2 핵심 요약}
이 강의에서는 연방준비제도 대차대조표와 그 구성 요소를 논의했습니다.
\begin{itemize}
    \item \textbf{주요 자산:} 연방준비제도가 보유한 주요 자산은 정부 증권입니다.
    \item \textbf{주요 부채:} 주요 부채는 본원 통화입니다. 이는 유통 통화와 은행 준비금으로 구성됩니다.
    \item \textbf{규모 변화:} 대차대조표 규모는 2008년과 2020년 사이에 10배 증가했습니다. 이때 연준은 경제를 지원하기 위해 통화정책 도구를 사용했습니다. 연준은 본원 통화, 특히 금융 시스템 내 은행 준비금의 양을 늘림으로써 이러한 지원에 대한 비용을 지불했습니다.
\end{itemize}
\end{summarybox}

\newpage
\section{중앙은행의 통화정책 도구와 비전통적 접근}

\begin{summarybox}
  \textbf{문서 전체 핵심 요약}
  이 문서는 중앙은행의 핵심 통화정책 도구인 공개시장운영과 재할인창구를 심층적으로 다룹니다.
  특히 2008년 금융위기 이후 등장한 비전통적 통화정책인 포워드 가이던스, 대규모 자산 매입(양적 완화)을 설명합니다.
  또한, 풍부한 은행 준비금 환경에서 금리 통제를 위한 새로운 도구인 지급준비금 이자(IORB)와 역환매조건부채권(ON RRP)의 역할과 작동 방식을 상세히 분석합니다.
  각 정책이 경제에 미치는 영향과 중앙은행의 목표 달성 과정을 이해하는 것이 이 문서의 주요 목표입니다.
\end{summarybox}

\newpage
\section{용어 정리}

\begin{table}[h!]
\centering
\caption{주요 통화정책 용어 및 설명}
\label{tab:glossary}
\begin{adjustbox}{width=\textwidth,center}
\begin{tabular}{lllc}
\toprule
\textbf{용어} & \textbf{쉬운 설명} & \textbf{원어} & \textbf{비고} \\
\midrule
공개시장운영 & 중앙은행이 국채 등을 사고팔아 시중의 돈(준비금)의 양을 조절하는 것. & Open Market Operations (OMO) & 가장 기본적인 통화정책 수단 \\
연방기금금리 & 은행들이 서로 하룻밤 동안 돈을 빌려주고 빌릴 때 적용되는 금리. & Federal Funds Rate & 미국 통화정책의 주요 목표 금리 \\
재할인창구 & 중앙은행이 은행들에게 단기 대출을 해주는 창구. & Discount Window & 은행의 단기 유동성 지원 및 금리 상한선 역할 \\
재할인율 & 재할인창구를 통해 은행이 중앙은행에서 돈을 빌릴 때 적용되는 금리. & Discount Rate & 연방기금금리의 상한선 역할을 함 \\
포워드 가이던스 & 중앙은행이 미래 통화정책 방향에 대해 시장에 미리 알려주는 소통 방식. & Forward Guidance & 불확실성 감소 및 장기 금리 영향 목적 \\
대규모 자산 매입 & 중앙은행이 장기 국채나 주택저당증권 등을 대량으로 매입하는 정책. & Large-Scale Asset Purchases (LSAPs) & 양적 완화(QE)라고도 불리며, 장기 금리 인하 목적 \\
양적 완화 & 대규모 자산 매입과 동일한 의미로 사용되며, 경제 활성화를 위한 비전통적 정책. & Quantitative Easing (QE) & 단기 금리가 0에 가까울 때 사용 \\
지급준비금 이자 & 중앙은행이 은행이 예치한 준비금에 대해 지급하는 이자. & Interest on Reserve Balances (IORB) & 풍부한 준비금 환경에서 금리 하한선 역할 \\
역환매조건부채권 & 중앙은행이 비은행 금융기관으로부터 증권을 담보로 돈을 빌리고 이자를 지급하는 것. & Overnight Reverse Repurchase Agreement (ON RRP) & 비은행 기관의 단기 금리 하한선 역할 \\
주요 딜러 & 중앙은행의 공개시장운영에 참여하는 대형 금융기관. & Primary Dealers & 주로 대형 은행들이 포함됨 \\
시스템 공개시장계정 & 뉴욕 연방준비은행이 연준의 공개시장운영을 관리하는 계정. & System Open Market Account (SOMA) & 연준의 자산 보유 현황을 보여줌 \\
\bottomrule
\end{tabular}
\end{adjustbox}
\end{table}

\newpage
\section{Lesson 2-2.3. 공개시장운영 (Open Market Operations)}

\begin{mybox}
  \textbf{핵심 요약}
  공개시장운영은 중앙은행이 국채 등의 증권을 사고팔아 시중의 은행 준비금 양을 조절하고, 이를 통해 연방기금금리(Federal Funds Rate)를 목표 수준으로 유도하는 핵심 통화정책 수단입니다.
\end{mybox}

공개시장운영(Open Market Operations, OMO)은 중앙은행의 통화정책 도구 중 가장 핵심적인 역할을 합니다. 만약 연방준비제도(Federal Reserve, Fed)가 금리를 인상하기로 결정한다면, 이는 공개시장운영을 통해 달성될 수 있습니다.

\subsection{1. 공개시장운영이란 무엇인가?}
공개시장운영은 중앙은행이 증권을 매각하거나 매입하는 행위를 말합니다. 이러한 행위가 실제로 어떻게 작동하는지, 그리고 연방기금시장(Federal Funds Market)에서 준비금(reserves)의 양에 직접적으로 어떤 영향을 미치는지 살펴보겠습니다. 궁극적으로 이러한 공개시장운영이 시장의 균형 이자율을 어떻게 결정하는지 이해하는 것이 중요합니다. 연방기금금리는 미국의 통화정책에서 주요 목표 이자율이라는 점을 기억해야 합니다.

\subsection{2. 통화정책의 목표 이자율: 연방기금금리}
통화정책은 목표 이자율을 선택하는 것을 포함합니다. 이 목표 이자율은 경제에 영향을 미쳐 연방준비제도가 목표를 달성할 수 있도록 돕습니다. 연방준비제도의 목표 이자율은 연방기금금리(Federal Funds Rate)라고 불리며, 이는 단기 이자율입니다. 사실상 이는 은행 간 시장(interbank market)의 이자율입니다. 미국 은행 간 시장은 연방기금시장이라고 불립니다. 모든 은행은 중앙은행에 준비금 계좌를 가지고 있으며, 이 준비금을 연방기금(federal funds)이라고 부릅니다. 공개시장운영은 이 연방기금금리에 영향을 미치도록 설계되었습니다.

\subsection{3. 공개시장운영의 작동 방식 (개념적 이해)}
연방준비제도는 이자율을 직접적으로 변경하지 않습니다. 대신, 공개시장운영을 통해 준비금의 공급을 변경합니다. 준비금의 공급과 수요의 교차점이 이 준비금의 가격을 결정하며, 이 가격이 바로 시장의 이자율, 즉 연방기금금리입니다.

\begin{examplebox}
  \textbf{예시: 금리 인하 목표 시}
  만약 연방준비제도가 이자율을 낮추고 싶다고 가정해 봅시다. 이를 달성하기 위해 연방준비제도는 준비금의 공급 곡선을 오른쪽으로 이동시킬 수 있습니다. 새로운 균형에서는 더 많은 연방기금이 더 낮은 이자율로 거래될 것입니다.
\end{examplebox}

\subsection{4. 준비금 조작을 통한 금리 통제}
공개시장운영은 연방준비제도가 시스템 내 준비금의 양을 조작할 수 있도록 합니다. 연준은 어떻게 이를 수행할까요?

\begin{examplebox}
  \textbf{예시: 은행 준비금 증가 목표 시}
  연준이 은행 준비금을 늘리고 싶다고 가정해 봅시다. 연준은 은행으로부터 증권을 매입하고 그 대가로 은행에 더 많은 준비금을 지급할 수 있습니다. 은행은 이미 이 증권을 소유하고 있었을 것이며, 연준은 공정한 가격을 지불해야 합니다.
\end{examplebox}

다시 한번 정리하면, 연방준비제도가 연방기금금리를 낮추고 싶다면, 은행으로부터 증권을 매입할 수 있습니다. 이러한 매입은 준비금으로 지급됩니다. 따라서 연방기금시장에서 전체 가용 자금은 증가할 것입니다. 이는 공급 곡선을 오른쪽으로 이동시킵니다. 새로운 균형에서는 더 많은 연방기금과 더 낮은 이자율이 형성됩니다.

\subsection{5. 공개시장운영의 실제 절차}
실제로 공개시장운영이 어떻게 작동하는지 몇 단계를 통해 살펴보겠습니다.

\begin{enumerate}
    \item \textbf{1단계: 목표 이자율 결정} \\
    연방준비제도 공개시장위원회(Federal Open Market Committee, FOMC)는 목표 이자율을 결정합니다.
    \item \textbf{2단계: 운영 위임} \\
    실제 증권 매입 및 매각은 뉴욕 연방준비은행(Federal Reserve Bank of New York)에 위임됩니다. 이들은 연방준비제도 시스템 공개시장계정(System Open Market Account, SOMA)을 관리합니다.
    \item \textbf{3단계: 증권 매입} \\
    뉴욕 연방준비은행의 트레이딩 데스크는 소위 주요 딜러(primary dealers)로부터 증권을 매입합니다. 이 딜러들은 연방준비제도에 준비금 계좌를 가지고 있으며, 가장 큰 은행들이 포함됩니다. 주로 미국 국채(US treasuries)가 매입됩니다.
    \item \textbf{4단계: 연준 대차대조표 확장} \\
    트레이딩 데스크가 매입을 완료하면 연방준비제도의 대차대조표가 확장됩니다. 자산 측면에서는 매입한 증권이 추가되고, 부채 측면에서는 준비금이 증가합니다. 이 준비금은 연방준비제도가 새로 창출하는 것입니다. 연준은 통화 기반(monetary base)을 통제하므로 이러한 준비금을 자유롭게 창출할 수 있습니다.
\end{enumerate}

\begin{examplebox}
  \textbf{예시: 연방준비제도의 공개시장 매입}
  2021년 7월 19일, 연준은 액면가 14억 달러의 미국 국채를 매입했습니다. 이 증권들은 15년에서 22년 사이의 잔여 만기를 가지고 있었습니다.
\end{examplebox}

\begin{figure}[h!]
  \centering
  % 실제 이미지가 없으므로 텍스트로 설명합니다.
  \begin{tcolorbox}[mybox, title=\textbf{시스템 공개시장계정 포트폴리오 (2021년 7월 19일)}]
    이 시점의 연방준비제도 시스템의 국내 증권 보유 현황을 보여주는 스냅샷입니다.
    모든 증권은 미국 정부 부채 또는 미국 정부가 보증하는 부채였습니다.
    포트폴리오는 이전 주 대비 약 1,000억 달러 증가했습니다.
  \end{tcolorbox}
  \caption{시스템 공개시장계정 포트폴리오 (2021년 7월 19일) 설명}
  \label{fig:soma_portfolio}
\end{figure}

\subsection{6. 2008년 금융위기 이후의 변화}
2008년 금융위기는 공개시장운영 방식에 전환점을 가져왔습니다. 연준은 2007년 중반 5.25\%였던 연방기금금리 목표를 2008년 12월 0\%에서 0.25\% 범위로 낮췄습니다. 단기 이자율이 거의 0에 가까워지자, 단순히 더 많은 증권을 매입하여 준비금을 늘리는 방식은 더 이상 효과적이지 않았습니다. 어려움을 겪는 경제를 자극하기 위한 새로운 접근 방식이 필요했습니다.

금융위기 이후, 은행 준비금의 총 공급량은 그 어느 때보다 많아졌습니다. 이는 연방준비제도의 대규모 개입의 직접적인 결과였습니다. 이는 전통적인 공개시장운영이 단기 이자율을 통제하는 데 덜 효과적일 것임을 의미했습니다. 연준은 이러한 문제에 대처하기 위한 새로운 통화정책 도구가 필요했습니다. 이 문제는 별도의 강의에서 논의될 것입니다.

\begin{summarybox}
  \textbf{핵심 정리: 공개시장운영}
  \begin{itemize}
    \item 공개시장운영은 준비금 공급에 영향을 미치며, 이는 연방준비제도의 목표 이자율인 연방기금금리를 결정합니다.
    \item 2008년 금융위기 이전에는 공개시장운영이 주요 통화정책 도구였습니다.
    \item 2008년 이후, 준비금 공급이 증가하면서 전통적인 공개시장운영의 중요성은 감소했습니다.
  \end{itemize}
\end{summarybox}

\newpage
\section{Lesson 2-2.4. 재할인창구 (The Discount Window)}

\begin{mybox}
  \textbf{핵심 요약}
  재할인창구는 중앙은행이 은행들에게 단기 현금 대출을 제공하는 전통적인 통화정책 도구입니다. 이는 은행의 유동성 관리를 돕고, 재할인율을 통해 연방기금금리의 상한선을 설정하여 통화정책에 영향을 미칩니다.
\end{mybox}

이 강의에서는 연방준비제도의 재할인창구(Discount Window)에 대해 논의합니다. 연방준비제도가 재할인창구를 사용하여 은행에 다양한 단기 현금 대출을 제공하는 방법을 살펴보겠습니다. 주요 신용(primary credit), 보조 신용(secondary credit), 계절 신용(seasonal credit)을 검토하고, 이러한 대출의 특징과 존재 이유를 이해할 것입니다. 그런 다음 재할인창구 대출과 통화정책 간의 연관성을 파악하고, 재할인율(discount rate)이 연방준비제도의 목표 이자율에 어떻게 영향을 미치는지 이해할 것입니다.

\subsection{1. 재할인 대출의 기본 개념}
재할인 대출(Discount lending)은 은행업의 오랜 전통입니다. 기본 아이디어는 대출 기관이 먼저 필요한 이자와 수수료 금액을 계산한다는 것입니다. 그런 다음 대출 기관은 이 금액을 대출의 액면가에서 할인합니다. 차용인은 오늘 할인된 대출 금액을 받습니다. 이 대출 금액은 이자와 수수료를 미리 반영하여 이미 줄어든 상태이지만, 차용인은 만기 시 전체 액면가를 상환해야 합니다.

\begin{examplebox}
  \textbf{예시: 재할인 대출 계산}
  액면가 100달러의 재할인 대출을 받고 싶다고 가정해 봅시다. 대출 만기는 1년입니다. 대출 기관이 10\%의 재할인율을 계산했다고 가정해 봅시다. 이는 대출 이자와 수수료를 포함합니다. 오늘 당신이 받는 대출 금액은 90달러입니다. 이는 액면가 100달러에서 이자와 수수료 10달러를 뺀 금액입니다. 하지만 1년 후에는 전체 100달러를 상환해야 합니다.
\end{examplebox}

\subsection{2. 재할인 대출의 유용성}
이러한 유형의 대출은 왜 유용할까요? 대출 비용이 선불로 처리되기 때문입니다. 이는 일련의 월별 이자 지급 대신 한 번의 상환만 필요하다는 것을 의미합니다. 이러한 이유로 재할인 대출은 종종 단기 자금 조달 필요성을 해결하는 데 도움이 됩니다.

\subsection{3. 연방준비제도의 재할인창구}
1913년부터 연방준비제도는 은행에 재할인 대출을 제공해 왔습니다. 연준은 왜 이런 일을 할까요? 이러한 대출은 은행에 단기간에 현금을 제공하여 은행 시스템에 충분한 현금이 있도록 보장합니다. 극단적인 경우, 이는 뱅크런(bank runs)과 은행 공황(banking panics)을 방지할 수 있습니다. 재할인 대출은 각 지역 연방준비은행의 특별 창구인 재할인창구에서 제공되었으며, 그래서 재할인창구 대출(discount window lending)이라는 이름이 붙었습니다. 오늘날 연준은 적격 담보로 담보된 단기 현금 대출을 제공하며, 이를 담보 재할인 대출(secured discount loans)이라고 합니다. 그러나 재할인창구라는 용어는 여전히 연방준비제도가 은행에 직접 대출할 때 사용됩니다.

\subsection{4. 재할인창구 대출의 종류: 주요 신용 (Primary Credit)}
연방준비제도에서 은행은 현재 여러 유형의 재할인 대출에 접근할 수 있습니다. 우리는 가장 중요한 범주인 주요 신용(primary credit)에 초점을 맞출 것입니다. 주요 신용의 주된 목적은 다른 자금 조달이 불가능한 경우 현금을 제공하는 것입니다. 주요 신용 재할인 대출은 자본이 충분하고 건전한 은행에 제공됩니다. 이들은 최대 만기 90일의 단기 대출입니다.

\begin{mybox}
  \textbf{주요 신용의 특징}
  \begin{itemize}
    \item \textbf{이자율:} 주요 신용 대출의 이자율은 목표 연방기금금리에 상대적으로 설정됩니다. 이는 일반적으로 다른 형태의 단기 신용 이자율보다 높습니다.
    \item \textbf{목적:} 이는 은행이 민간 현금 시장에서 차입하도록 유도하기 위함입니다. 연준은 최후의 수단(last resort)이 되어야 합니다.
    \item \textbf{담보:} 이들은 담보 대출입니다. 연방준비제도는 상업 및 산업 대출, 소비자 대출, 주거 및 상업용 부동산 대출, 미국 국채를 포함한 광범위한 담보를 허용합니다.
    \item \textbf{마진:} 일부 담보는 다른 담보보다 위험합니다. 위험을 반영하기 위해 연준은 은행이 각 유형의 담보에 대해 얼마나 빌릴 수 있는지를 나타내는 마진(margins)을 적용합니다.
  \end{itemize}
\end{mybox}

\begin{figure}[h!]
  \centering
  % 실제 이미지가 없으므로 텍스트로 설명합니다.
  \begin{tcolorbox}[mybox, title=\textbf{마진 테이블 예시}]
    이 마진은 시장 가치의 백분율로 표시됩니다.
    \begin{itemize}
      \item 잔여 만기 1년의 미국 국채는 매우 안전합니다. 은행은 1달러당 99센트를 빌릴 수 있습니다.
      \item 대조적으로, 잔여 만기 1년의 상업용 부동산 대출은 여전히 상당히 위험합니다. 은행은 대출 가치의 35\%에서 95\% 사이만 빌릴 수 있습니다.
    \end{itemize}
  \end{tcolorbox}
  \caption{담보 유형별 마진 예시 설명}
  \label{fig:margin_table}
\end{figure}

\subsection{5. 재할인창구와 통화정책의 연관성}
재할인창구는 연방준비제도가 경제의 단기 이자율을 통제하는 데 어떻게 도움이 될까요? 두 가지를 염두에 두어야 합니다.
\begin{enumerate}
    \item \textbf{대체재 관계:} 연준으로부터의 재할인 대출과 다른 은행으로부터의 은행 간 대출은 밀접한 대체재입니다. 둘 다 단기 대출입니다. 모든 건전한 은행은 은행 간 시장이나 재할인창구에 접근할 수 있습니다.
    \item \textbf{재할인율의 상한선 역할:} 재할인율은 연방기금 목표 범위의 상한선 근처에 설정된다는 점을 기억하십시오.
\end{enumerate}

그렇다면 연방기금의 실현 이자율이 재할인율을 초과하면 어떻게 될까요? 은행이 재할인창구에서 빌리는 것이 더 저렴할 것입니다. 은행 간 대출에서 멀어지는 이러한 수요 이동은 연방기금금리에 하향 압력을 가할 것입니다. 따라서 재할인율은 중앙은행이 설정한 연방기금 목표의 상한선을 강제합니다. 이는 실현된 연방기금금리가 너무 높아지지 않도록 보장합니다.

\begin{figure}[h!]
  \centering
  % 실제 이미지가 없으므로 텍스트로 설명합니다.
  \begin{tcolorbox}[mybox, title=\textbf{주요 신용 금리와 연방기금금리 비교}]
    데이터에서 이를 확인할 수 있습니다. 2003년부터 주요 신용 금리는 항상 연방기금금리보다 높게 유지되었습니다.
  \end{tcolorbox}
  \caption{주요 신용 금리와 연방기금금리 추이 설명}
  \label{fig:primary_credit_vs_fed_funds}
\end{figure}

\begin{summarybox}
  \textbf{핵심 정리: 재할인창구}
  \begin{itemize}
    \item 재할인창구 대출은 은행에 대한 단기 담보 대출입니다. 이는 은행이 단기 현금 필요성을 관리하는 데 도움이 됩니다.
    \item 재할인 대출은 은행 간 시장에서 연방기금을 차입하는 것의 대체재입니다.
    \item 결과적으로 재할인율은 연방기금금리의 상한선을 제공합니다.
  \end{itemize}
\end{summarybox}

\newpage
\section{Lesson 2-3. 비전통적 통화정책 (Non-traditional Monetary Policy)}

\subsection{Lesson 2-3.1. 포워드 가이던스 (Forward Guidance)}

\begin{mybox}
  \textbf{핵심 요약}
  포워드 가이던스는 중앙은행이 미래 통화정책 의도를 시장에 소통하여 시장의 기대를 형성하고 불확실성을 줄이는 비전통적 통화정책 도구입니다. 이는 단기 금리가 0에 가까워졌을 때 장기 금리에 영향을 미치기 위해 도입되었습니다.
\end{mybox}

이 강의에서는 새로운 통화정책 도구인 포워드 가이던스(Forward Guidance)를 살펴보겠습니다. 포워드 가이던스는 연방준비제도가 미래 통화정책 의도를 소통할 수 있도록 합니다. 이는 오늘 취해진 조치보다는 미래의 조치에 대한 소통입니다. 이 소통은 미래 인플레이션이나 경제 성장에 관한 것일 수 있으며, 종종 연준이 미래에 이자율을 어떻게 조정할 계획인지 포함합니다. 목표는 시장 기대를 형성하고 불확실성을 줄이는 것입니다. 포워드 가이던스의 여러 예시를 논의하고, 연준이 왜 포워드 가이던스를 사용하는지, 그리고 그것이 금융 상황에 어떻게 영향을 미치는지 이해할 것입니다.

\subsection{1. 포워드 가이던스의 등장 배경}
포워드 가이던스는 비전통적 통화정책 도구입니다. 전통적인 통화정책은 단기 이자율에 작용하며, 실제로는 공개시장운영과 재할인창구 대출을 의미합니다. 포워드 가이던스는 비교적 새로운 개념으로, 2009년 연방기금금리가 0에 가까워졌을 때 등장했습니다. 이 시점에서는 단기 이자율을 목표로 하는 것이 더 이상 선택 사항이 아니었습니다. 이 문제를 극복하기 위해 연방공개시장위원회(FOMC)는 미래에 통화정책을 어떻게 조정할 의도인지 소통하기로 결정했습니다. 연준은 장기 이자율을 목표로 삼으려 했습니다.

\subsection{2. 포워드 가이던스의 형태와 목표}
포워드 가이던스는 FOMC의 공식적인 소통입니다. 이는 보도 자료, 연설, 의회 증언 등의 형태로 나타날 수 있습니다. 이들은 일반적으로 대중에게 매우 신중하게 작성된 성명서로, 혼란을 피하는 것이 목표입니다. 포워드 가이던스는 대중에게 통화정책 의도를 알립니다. 이는 향후 몇 달 또는 몇 년 동안의 조치일 수 있으며, 경제가 어떻게 진화하는지에 따라 달라질 수 있습니다. 목표는 장기 이자율을 움직이는 것일 수 있습니다. 포워드 가이던스는 또한 불확실성을 줄일 수 있습니다.

\subsection{3. 초기 FOMC 성명서의 특징 (2009년 이전)}
연준이 어떻게 소통했는지에 대한 초기 예시를 먼저 살펴보겠습니다. 1994년부터 FOMC는 가장 최근의 통화정책 회의 결정 사항을 발표하는 성명서를 발행했습니다. 2009년 이전에는 이러한 성명서가 간결했습니다.

\begin{examplebox}
  \textbf{예시: 2004년 3월 16일 성명서}
  이 성명서는 이자율에 대한 한 문장, 현재 경제 상황을 요약한 한 단락, 그리고 잠재적 위험에 대한 한 단락으로 구성되어 있습니다. 또한 어떤 FOMC 위원들이 조치에 찬성했고, 반대한 위원이 있다면 누구인지도 명시합니다. 마지막으로, 미래 통화정책 조치에 대한 어떠한 언급도 없다는 점에 주목하십시오.
\end{examplebox}

\begin{tcolorbox}[mybox, title=\textbf{성명서 비교를 통한 시장 분석 (2004년 1월 vs 3월)}]
  2004년 1월 28일 성명서와 비교해 봅시다. 두 번째 단락에서 경제에 대한 설명이 약간 변경되었습니다. 1월에는 확장이 위험했지만, 3월에는 "견고한 속도의 확장"에 불과했습니다. 이는 성장률이 여전히 긍정적이지만 감소했음을 나타냅니다. 마찬가지로 노동 시장에 대한 설명도 덜 우호적입니다. 마지막 단락은 "that"이라는 단어의 삽입을 제외하고는 동일합니다. 이러한 언어의 변화는 대중에게 면밀히 관찰됩니다.
  
  \textbf{시사점:} 이 두 성명서를 비교하여 무엇을 알 수 있을까요? 6주 동안 연방준비제도는 경제 상황에 대한 의견을 변경했습니다. 1월 성명서는 상당히 낙관적이었지만, 3월 성명서는 덜 낙관적입니다. 이것이 통화정책에 무엇을 의미할까요? 아마도 연방준비제도는 다음 회의에서 이자율을 인상하지 않을 것입니다. 이는 사실입니다. 비록 미래 통화정책에 대한 명확한 언급은 포함되어 있지 않지만 말입니다.
\end{tcolorbox}

\subsection{4. 명시적 포워드 가이던스의 등장 (2009년 이후)}
이제 더 최근의 FOMC 성명서를 살펴보겠습니다. 여기에는 명시적인 포워드 가이던스가 포함되어 있습니다.

\begin{examplebox}
  \textbf{예시: 2009년 3월 18일 성명서}
  이 시점에서 연방기금금리는 0에 가까웠습니다. FOMC는 무엇이라고 말했을까요?
  "위원회는 연방기금금리 목표 범위를 0에서 0.25\%로 유지할 것이며, 경제 상황이 \textbf{상당히 오랜 기간 동안(for an extended period)} 연방기금금리의 예외적으로 낮은 수준을 정당화할 가능성이 있다고 예상합니다."
  
  이것이 바로 포워드 가이던스입니다. 이 성명서는 연방기금금리가 "상당히 오랜 기간 동안" 예외적으로 낮은 수준, 즉 0으로 유지될 것임을 시사합니다. 시간 프레임은 정확히 명확하지 않습니다. 대중은 이를 어떻게 해석해야 할까요? 금융 시장 분석가들은 아마도 연방기금금리가 최소 1년, 그리고 훨씬 더 오랫동안 0으로 유지될 것으로 예상했을 것입니다.
\end{examplebox}

\begin{examplebox}
  \textbf{예시: 2012년 12월 12일 성명서 (조건부 포워드 가이던스)}
  이제 2012년 12월 12일 성명서를 고려해 봅시다. 이 시점에서 연방기금금리는 여전히 0에 가까웠지만, 금리 인상이 가능성이 있었습니다. 이 성명서는 훨씬 더 길다는 것을 알 수 있습니다. FOMC는 경제에 대해 더 자세히 설명합니다. 그들은 가계 및 기업 투자와 같은 특정 부문을 언급합니다.
  
  그러나 여기에서 FOMC는 새로운 것을 합니다. 5번째 단락에서 다음과 같이 명시합니다.
  "위원회는 연방기금금리 목표 범위를 0에서 0.25\%로 유지하기로 결정했으며, 현재 이 예외적으로 낮은 연방기금금리 범위가 적절할 것으로 예상합니다. \textbf{적어도 실업률이 6.5\%를 초과하는 한, 1~2년 후 인플레이션이 위원회의 2\% 장기 목표보다 0.5\% 포인트 이상 높지 않을 것으로 예상되는 한, 그리고 장기 인플레이션 기대치가 잘 고정되어 있는 한 말입니다.} 이 문장은 상당히 길고 복잡합니다."
\end{examplebox}

\begin{tcolorbox}[mybox, title=\textbf{조건부 포워드 가이던스 해설}]
  이것을 풀어봅시다. FOMC는 다음과 같은 정확한 조건 하에서 목표 연방기금금리를 0으로 유지할 것이라고 말합니다.
  \begin{enumerate}
    \item 실업률이 6.5\%를 초과하는 한.
    \item 1~2년 후 예상 인플레이션이 2.5\%를 초과하지 않는 한.
    \item 인플레이션 기대치가 잘 고정되어 있는 한 (즉, 모든 사람이 장기적으로 2\% 인플레이션을 여전히 예상하는 한).
  \end{enumerate}
  
  다시 말해, FOMC는 이러한 정확한 조건 하에서 이자율을 낮게 유지할 것입니다. 이러한 유형의 포워드 가이던스는 계획의 확실성을 제공합니다. 이는 "상당히 오랜 기간"이 무엇을 의미하는지 명확히 합니다.
\end{tcolorbox}

\subsection{5. FOMC 회의록 및 경제 전망 요약 (SEP)}
FOMC는 또한 금리 설정 회의록을 공개합니다. 이들은 회의에서 논의된 내용을 그대로 기록한 것입니다. 회의록은 각 회의 후 3주 후에 공개됩니다. 회의록은 논의된 정책 관련 주제, 참가자들이 표명한 견해, 전망에 대한 위험과 불확실성, 위원회 결정의 이유에 대한 추가 통찰력을 제공합니다. 물론 금융 시장 분석가들은 미래 통화정책에 대한 통찰력을 얻기 위해 이 회의록을 면밀히 검토합니다.

2007년부터 FOMC는 경제 전망 요약(Summary of Economic Projections, SEP)도 발행하고 있습니다. 이는 매년 4회 발행됩니다. 이 전망은 FOMC 참가자들의 미래 경제 결과에 대한 평가를 제공합니다. 여기에는 GDP, 실업률, 인플레이션, 이자율이 포함됩니다. 이러한 결과는 중기 및 장기적으로 고려됩니다. 각 FOMC 참가자가 평가를 제공하므로 대중은 다양한 예측 범위를 볼 수 있습니다.

\begin{figure}[h!]
  \centering
  % 실제 이미지가 없으므로 텍스트로 설명합니다.
  \begin{tcolorbox}[mybox, title=\textbf{2021년 3월 경제 전망 요약 (실질 GDP 변화)}]
    2021년 3월 17일 경제 전망 요약을 살펴보겠습니다. 당시 경제는 COVID-19 팬데믹으로 인해 어려움을 겪고 있었습니다.
    
    \textbf{실질 GDP 변화 차트:} 인플레이션을 고려한 GDP 성장률입니다. 2021년 예측은 5\%에서 7.3\% 사이로 다양하다는 것을 알 수 있습니다. 단기 예측의 차이는 회복에 대한 불확실성을 반영합니다. 장기적으로는 참가자들의 추정치가 서로 매우 가깝습니다. 장기적으로는 의견 불일치가 거의 없습니다.
  \end{tcolorbox}
  \caption{2021년 3월 실질 GDP 변화 전망 설명}
  \label{fig:sep_gdp}
\end{figure}

\begin{figure}[h!]
  \centering
  % 실제 이미지가 없으므로 텍스트로 설명합니다.
  \begin{tcolorbox}[mybox, title=\textbf{2021년 3월 경제 전망 요약 (닷 플롯)}]
    2021년 3월 경제 전망 요약에서 소위 닷 플롯(dot plots)도 있습니다. 이 차트는 FOMC가 연방기금금리가 어떻게 될 것이라고 생각하는지를 요약합니다. 금융 분석가들은 이 포워드 가이던스에 세심한 주의를 기울입니다.
    
    \textbf{닷 플롯 분석:} 당시 이자율은 0에 가까웠습니다. 질문은 금리가 얼마나 빨리 인상될 것인가였습니다.
    \begin{itemize}
      \item 2021년에는 어떤 참가자도 0에서 인상을 예상하지 않았습니다.
      \item 2022년에는 4명의 참가자가 경제가 연방기금금리를 최소 한 번 인상할 만큼 충분히 회복될 것이라고 생각했습니다.
      \item 2023년에는 거의 절반의 참가자가 연방기금금리가 인상될 것으로 보았습니다.
      \item 장기적으로 참가자들은 연방기금금리가 2\%에서 3\% 사이가 될 것으로 예상합니다.
    \end{itemize}
    물론 이러한 예측은 수정될 수 있습니다. 그러나 이들은 FOMC 참가자들의 사고에 대한 귀중한 통찰력을 제공합니다. 이러한 종류의 포워드 가이던스는 금융 시장 참가자들에게 더 많은 계획의 확실성을 제공합니다.
  \end{tcolorbox}
  \caption{2021년 3월 닷 플롯 설명}
  \label{fig:sep_dot_plot}
\end{figure}

\begin{summarybox}
  \textbf{핵심 정리: 포워드 가이던스}
  \begin{itemize}
    \item 포워드 가이던스는 중앙은행이 경제 및 통화정책이 어떻게 진화할 것으로 예상하는지 신호를 보낼 수 있는 통화정책 도구입니다.
    \item 대중, 특히 금융 시장 참가자들은 이에 세심한 주의를 기울입니다.
    \item 일상적인 포워드 가이던스는 단기 이자율이 매우 낮았던 2009년에 등장했습니다.
    \item 동시에 연방준비제도는 장기 이자율에 압력을 가하기 위해 포워드 가이던스를 도입했습니다.
  \end{itemize}
\end{summarybox}

\newpage
\subsection{Lesson 2-3.2. 대규모 자산 매입 (Large-Scale Asset Purchases)}

\begin{mybox}
  \textbf{핵심 요약}
  대규모 자산 매입(LSAPs), 또는 양적 완화(QE)는 단기 금리가 0에 가까울 때 장기 금리에 직접적으로 영향을 미쳐 경제를 자극하기 위해 사용되는 비전통적 통화정책입니다. 중앙은행이 장기 채권을 대량으로 매입하여 채권 가격을 올리고 수익률을 낮춥니다.
\end{mybox}

이 강의에서는 대규모 자산 매입(Large-Scale Asset Purchases, LSAPs)에 대해 다룹니다. LSAPs는 종종 양적 완화(Quantitative Easing, QE)라고 불립니다. 이들은 장기 이자율에 직접적으로 영향을 미치도록 설계된 비전통적 통화정책의 한 형태입니다. 연방준비제도는 2008년 금융위기 동안 이 도구를 처음 사용했습니다. 당시 연준은 단기 이자율이 0이었음에도 불구하고 경제를 자극하려고 노력했습니다. 이 세션에서는 LSAPs가 어떻게 작동하는지, 이러한 프로그램이 어떻게 설계되었는지, 장기 이자율에 어떻게 영향을 미쳤는지, 그리고 경제 전반에 미친 광범위한 영향을 살펴보겠습니다.

\subsection{1. 대규모 자산 매입의 필요성}
전통적인 통화정책은 단기 이자율, 일반적으로 연방기금금리 또는 재할인율에만 작용합니다. 연준이 단기 이자율을 인하하면 장기 이자율도 낮아질 수 있습니다. 그러나 장기 이자율은 덜 움직입니다. 2008년 후반, 미국 경제는 심각한 경기 침체에 진입하고 있었습니다. 연방준비제도는 경제를 자극하고 싶었지만, 연방기금금리는 이미 0으로 인하된 상태였습니다. 통화정책에 대한 새로운 접근 방식이 필요했습니다.

\subsection{2. 중앙은행의 대응: QE 프로그램}
중앙은행은 이 문제에 어떻게 대응했을까요? 연방준비제도는 대량의 장기 채권을 매입하기 시작했습니다. 이들이 소위 QE 프로그램입니다. QE1, QE2, QE3, 그리고 오퍼레이션 트위스트(Operation Twist)의 네 가지 프로그램이 있었습니다.

\begin{mybox}
  \textbf{QE 프로그램의 작동 원리}
  \begin{itemize}
    \item \textbf{목표:} 이러한 매입은 장기 채권에 대한 대규모 수요 충격처럼 작용했습니다.
    \item \textbf{가격 상승:} 초과 수요는 가격을 상승시킵니다.
    \item \textbf{수익률 하락:} 가격이 높다는 것은 채권 수익률이 낮다는 것을 의미합니다. 채권 가격과 채권 수익률 사이에는 역관계가 있다는 것을 기억하십시오.
    \item \textbf{경제 자극:} 따라서 이러한 매입은 장기 채권 수익률을 직접적으로 낮추는 것을 목표로 합니다. 장기 채권 수익률이 낮아지면 장기 투자 기간에 대한 차입 비용이 줄어듭니다. 이는 투자와 경제 활동을 활성화해야 합니다.
  \end{itemize}
\end{mybox}

\subsection{3. QE1 프로그램의 규모와 효과}
QE1 프로그램의 맥락에서 이러한 프로그램을 살펴보고, 그 규모와 효과를 이해해 봅시다. QE1은 2008년 11월에 시작되어 2009년 3월에 확대되었습니다. 2010년 3월에 종료될 때까지 연준은 1조 2,500억 달러의 주택저당증권(mortgage-backed securities, MBS), 1,750억 달러의 GSE(Fannie Mae 및 Freddie Mac) 발행 부채, 그리고 3,000억 달러의 미국 국채를 매입했습니다. 이러한 매입은 당시 시장 가치의 20\% 이상을 차지했습니다. 이는 중앙은행의 대규모 개입입니다.

첫 번째 양적 완화는 성공적이었을까요? 장기 수익률을 살펴보겠습니다.

\begin{figure}[h!]
  \centering
  % 실제 이미지가 없으므로 텍스트로 설명합니다.
  \begin{tcolorbox}[mybox, title=\textbf{QE1의 장기 국채 수익률 영향}]
    \textbf{장기 국채:} 2년 만기 국채 수익률 차트를 보면, 프로그램이 시행된 후 2008년 10월 후반 1.6\%에서 2009년 4월 1\%로 떨어졌습니다.
  \end{tcolorbox}
  \caption{QE1 시행 후 2년 만기 국채 수익률 변화 설명}
  \label{fig:qe1_treasury_yield}
\end{figure}

\begin{figure}[h!]
  \centering
  % 실제 이미지가 없으므로 텍스트로 설명합니다.
  \begin{tcolorbox}[mybox, title=\textbf{QE1의 회사채 수익률 영향}]
    \textbf{회사채:} BAA 등급 회사채 수익률을 살펴보겠습니다. 이는 자본 접근성이 제한적인 기업의 차입 비용입니다. 차트에서 볼 수 있듯이, QE1이 시작되면서 수익률은 최고 9.5\%에서 약 8\%로 떨어졌습니다. 이는 기업의 차입 비용을 낮췄습니다.
  \end{tcolorbox}
  \caption{QE1 시행 후 BAA 등급 회사채 수익률 변화 설명}
  \label{fig:qe1_corporate_bond_yield}
\end{figure}

\begin{figure}[h!]
  \centering
  % 실제 이미지가 없으므로 텍스트로 설명합니다.
  \begin{tcolorbox}[mybox, title=\textbf{QE1의 주택담보대출 시장 영향}]
    \textbf{주택담보대출 시장:} 연준은 주택저당증권(MBS)을 대량으로 매입했습니다. MBS 수익률과 주택담보대출 금리 사이에는 밀접한 관계가 있습니다. MBS 수익률이 하락하면 일반적으로 주택담보대출 금리도 하락합니다. 차트에서 QE1 프로그램이 시작되었을 때 30년 고정 주택담보대출 금리가 2008년 10월 후반 거의 6.5\%에서 2009년 4월 초 4.85\%로 떨어지기 시작한 것을 볼 수 있습니다. 이는 상당한 효과를 가져왔습니다. 주택담보대출 금리 인하는 주택 구매를 장려했습니다. 이러한 개입은 주택 시장을 지원하는 데 도움이 되었을 수 있습니다.
  \end{tcolorbox}
  \caption{QE1 시행 후 30년 고정 주택담보대출 금리 변화 설명}
  \label{fig:qe1_mortgage_rate}
\end{figure}

\subsection{4. 다른 QE 프로그램과 연준 대차대조표의 변화}
다른 양적 완화 프로그램들도 비슷한 성격을 가졌습니다. QE2는 2010년 10월에 발표되었고, QE3는 2012년 9월에 발표되었습니다. 오퍼레이션 트위스트(Operation Twist)는 연방준비제도가 단기 증권을 매각하고 그 수익금으로 장기 증권을 매입하는 것을 포함했습니다. 각 경우에 의도는 장기 이자율을 낮추는 것이었습니다. 이는 가계와 기업의 차입 비용을 줄입니다. 이러한 프로그램은 2014년 10월에 중단되었습니다. 이 시점에는 노동 시장이 회복되었습니다.

\begin{figure}[h!]
  \centering
  % 실제 이미지가 없으므로 텍스트로 설명합니다.
  \begin{tcolorbox}[mybox, title=\textbf{QE 프로그램이 연준 대차대조표에 미친 영향}]
    이 차트는 QE 프로그램이 연방준비제도의 대차대조표에 미친 영향을 보여줍니다. 이러한 대규모 자산 매입은 결국 연준의 자산으로 귀결되었습니다.
    
    2014년까지 연방준비제도는 거의 4조 달러에 달하는 증권을 보유했습니다. 이들은 주로 미국 국채와 MBS였습니다. 이는 2008년 금융위기 이전 기간에 비해 대차대조표의 극적인 확장입니다.
  \end{tcolorbox}
  \caption{QE 프로그램에 따른 연준 대차대조표 변화 설명}
  \label{fig:fed_balance_sheet_qe}
\end{figure}

\subsection{5. COVID-19 팬데믹 대응으로서의 LSAPs}
대규모 자산 매입은 COVID-19 팬데믹에 대한 연방준비제도의 대응의 일부이기도 했습니다. 2020년 3월, FOMC는 연방준비제도가 "원활한 시장 기능과 통화정책의 광범위한 금융 상황으로의 효과적인 전달을 지원하는 데 필요한 만큼" 국채와 기관 주택저당증권을 계속 매입할 것이라고 발표했습니다. 이 프로그램은 경제를 지원하기 위한 또 다른 개방형 약속이었습니다.

\begin{summarybox}
  \textbf{핵심 정리: 대규모 자산 매입 (LSAPs/QE)}
  \begin{itemize}
    \item 단기 이자율이 0\%일 때, 연방준비제도는 대규모 자산 매입(QE)을 사용하여 통화정책을 수행할 수 있습니다.
    \item 대규모 자산 매입은 장기 이자율을 목표로 할 수 있습니다.
    \item 대규모 자산 매입은 장기 이자율을 낮춤으로써 소비자와 기업의 신용 비용을 줄입니다.
  \end{itemize}
\end{summarybox}

\newpage
\subsection{Lesson 2-3.3. 풍부한 은행 준비금 환경에서의 통화정책 (Monetary Policy With Abundant Bank Reserves)}

\begin{mybox}
  \textbf{핵심 요약}
  2008년 금융위기 이후 대규모 자산 매입으로 인해 은행 준비금이 풍부해지면서, 전통적인 공개시장운영의 금리 통제 효과가 감소했습니다. 이에 연준은 지급준비금 이자(IORB)와 역환매조건부채권(ON RRP)이라는 새로운 도구를 도입하여 풍부한 준비금 환경에서도 단기 금리를 효과적으로 통제할 수 있게 되었습니다.
\end{mybox}

이 강의에서는 은행 준비금이 풍부할 때 통화정책을 어떻게 구현하는지 다룹니다. 2008년 이후, 그리고 대규모 자산 매입 프로그램 이후, 은행 시스템은 준비금으로 넘쳐났습니다. 이는 전통적인 공개시장운영이 이자율을 관리하는 데 덜 효과적이게 만들었습니다. 이 강의에서는 이러한 환경을 위해 만들어진 두 가지 새로운 통화정책 도구를 검토합니다. 첫째, 지급준비금 이자(Interest on Reserve Balances, IORB)와 둘째, 역환매조건부채권(Overnight Reverse Repurchase Agreement, ON RRP) 시설입니다. 이러한 도구는 은행 준비금 수준이 높을 때 연방준비제도가 단기 이자율을 통제하는 데 도움이 됩니다.

\subsection{1. 전통적인 금리 목표 방식의 한계}
연방준비제도는 어떻게 이자율을 목표로 삼을까요? 역사적으로 연방준비제도는 공개시장운영을 수행했습니다. 이는 연방기금시장에서 준비금 공급 수준을 변경하기 위해 증권을 사고파는 것을 포함합니다.

\begin{figure}[h!]
  \centering
  % 실제 이미지가 없으므로 텍스트로 설명합니다.
  \begin{tcolorbox}[mybox, title=\textbf{전통적인 연방기금시장 모델}]
    이 그림은 연방기금시장을 보여줍니다. 은행 준비금의 공급을 조정함으로써 연방준비제도는 연방기금금리를 통제합니다.
    \textbf{그래프 설명:} 수직축은 연방기금금리, 수평축은 준비금의 양을 나타냅니다. 준비금 공급 곡선은 수직선이며, 수요 곡선은 우하향합니다. 공급 곡선을 이동시켜 균형 금리를 조절합니다.
  \end{tcolorbox}
  \caption{전통적인 연방기금시장 모델 설명}
  \label{fig:traditional_fed_funds_market}
\end{figure}

그러다 2008년 금융위기가 닥쳤습니다. 연방기금금리는 0으로 인하되었고, 대규모 자산 매입이 이루어졌습니다. 이러한 매입은 은행 준비금의 급격한 증가로 자금이 조달되었습니다. 준비금은 2007년 140억 달러에서 2014년 2조 6천억 달러 이상으로 증가했습니다.

\begin{figure}[h!]
  \centering
  % 실제 이미지가 없으므로 텍스트로 설명합니다.
  \begin{tcolorbox}[mybox, title=\textbf{풍부한 준비금 환경에서의 연방기금시장}]
    우리는 다른 강의에서 이 차트를 검토했습니다. 연방기금시장의 균형에 미치는 영향을 볼 수 있습니다. 준비금 공급이 훨씬 더 많아졌고, 연방기금금리는 0입니다.
    \textbf{그래프 설명:} 준비금 공급 곡선이 극도로 오른쪽으로 이동하여 수요 곡선의 거의 평평한 부분과 교차합니다. 이로 인해 준비금의 양이 크게 증가했음에도 불구하고 연방기금금리는 0에 가깝게 유지됩니다.
  \end{tcolorbox}
  \caption{풍부한 준비금 환경에서의 연방기금시장 설명}
  \label{fig:abundant_reserves_fed_funds_market}
\end{figure}

\subsection{2. 풍부한 준비금 환경의 문제점}
이것이 왜 문제일까요? 결국 연방준비제도는 이자율을 인상하고 싶을 것입니다. 그러나 연방기금금리는 더 이상 은행 준비금의 작은 변화에 반응하지 않을 것입니다. 전통적인 공개시장운영만으로는 충분하지 않을 것입니다. 연방기금금리가 반응하려면 준비금 공급에 큰 변화가 필요했습니다. 이론적으로 연방준비제도는 대규모 자산 매각을 통해 준비금을 줄일 수 있었습니다. 매입이 아닌 자산 매각 말입니다. 그러나 이는 너무 비용이 많이 들 것이라는 우려가 있었습니다. 1조 달러의 증권을 매각하는 것은 금융 시장과 차용인에게 큰 의도치 않은 결과를 초래할 수 있었습니다. 이는 너무 위험하다고 판단되었습니다.

\subsection{3. 새로운 통화정책 도구의 도입}
대신 두 가지 대안이 채택되었습니다.
\begin{enumerate}
    \item \textbf{지급준비금 이자 (IORB):} 연방준비제도는 이제 지급준비금 잔액에 대해 이자를 지급합니다. 이 도구는 은행을 대상으로 합니다.
    \item \textbf{역환매조건부채권 시설 (ON RRP):} 연방준비제도는 역환매조건부채권 시설(Overnight Reverse Repurchase Agreement, ON RRP)을 만들었습니다. ON RRP는 머니마켓 뮤추얼 펀드와 같은 비은행 기관을 대상으로 합니다. 이러한 비은행 기관도 현금 시장의 활발한 참여자이므로 무시할 수 없습니다.
\end{enumerate}
이러한 도구들은 함께 연방준비제도가 준비금이 풍부할 때에도 단기 이자율을 통제할 수 있도록 합니다.

\subsection{4. 지급준비금 이자 (IORB)의 작동 방식}
어떻게 작동하는지 살펴보겠습니다. 지급준비금 잔액에 대한 이자부터 시작하겠습니다. 연방준비제도는 이제 지급준비금 잔액에 대해 이자를 지급합니다. 이들은 은행이 연준 계좌에 하룻밤 동안 보유하는 자금입니다. 이 자금은 은행 간 시장에서 다른 은행에 대출되지 않습니다. 이것이 통화정책 구현에 어떻게 도움이 될까요?

\begin{mybox}
  \textbf{IORB의 금리 하한선 역할}
  연방준비제도가 지급준비금 잔액에 대한 이자를 목표 연방기금금리와 동일하게 설정한다고 가정해 봅시다. 이제 은행들은 연준에 하룻밤 동안 투자하는 것과 은행 간 대출을 하는 것 사이에서 무관심해집니다. 어느 쪽이든 목표 연방기금금리를 벌게 될 것입니다. 이는 균형 상태에서 반드시 그래야 합니다. 만약 연방기금금리가 IORB보다 높다면, 은행들은 준비금을 은행 간 시장에 대출할 것입니다. 이러한 차익 거래 과정은 두 이자율이 같아질 때까지 계속되어야 합니다. 실제로 지급준비금 이자는 연방기금금리의 하한선을 제공합니다.
\end{mybox}

\begin{figure}[h!]
  \centering
  % 실제 이미지가 없으므로 텍스트로 설명합니다.
  \begin{tcolorbox}[mybox, title=\textbf{IORB가 연방기금금리에 미치는 영향}]
    여기 연방기금시장의 공급과 수요가 있습니다. 지급준비금 잔액에 이자를 지급함으로써 이는 연방기금금리에 하한선을 설정할 것입니다. 균형 상태에서 어떤 은행도 은행 간 시장에 대출하여 IORB보다 적게 받으려 하지 않을 것입니다.
    \textbf{그래프 설명:} 준비금 수요 곡선이 IORB 수준에서 평평해지는 구간을 보여줍니다. 이는 IORB가 연방기금금리의 효과적인 하한선으로 작용함을 나타냅니다.
  \end{tcolorbox}
  \caption{IORB가 연방기금금리에 미치는 영향 설명}
  \label{fig:iorb_fed_funds_market}
\end{figure}

이는 연방준비제도가 준비금 공급을 변경하지 않고도 연방기금금리를 인상할 수 있도록 합니다.

\subsection{5. 역환매조건부채권 시설 (ON RRP)의 작동 방식}
ON RRP 시설은 왜 필요할까요? 어떻게 작동할까요? 미국에서는 비은행 기관이 현금 시장에서 중요한 역할을 합니다. 예를 들어, 머니마켓 펀드(money market funds)가 있습니다. 이러한 대출 기관은 IORB에 접근할 수 없기 때문에 때로는 더 낮은 금리로 대출할 의향이 있을 수 있습니다. 이는 유휴 현금에 대해 아무것도 받지 않는 것보다 낫습니다. 이는 단기 이자율이 IORB 아래로 떨어질 수 있음을 의미합니다. 이러한 효과에 대처하기 위해 연방준비제도는 ON RRP 시설을 만들었습니다. 이는 이러한 비은행 기관을 위한 IORB와 동일합니다. 그들은 환매조건부채권(repurchase agreement)을 통해 연준에 하룻밤 동안 투자할 수 있습니다. 이 시설은 연준이 단기 이자율에 하한선을 설정할 수 있도록 보장합니다. 이러한 비은행 기관은 ON RRP 금리보다 적게 대출하려 하지 않을 것입니다.

\begin{figure}[h!]
  \centering
  % 실제 이미지가 없으므로 텍스트로 설명합니다.
  \begin{tcolorbox}[mybox, title=\textbf{ON RRP 금리와 연방기금금리 추이}]
    이 차트는 ON RRP 금리와 연방기금금리를 함께 보여줍니다. 보시다시피 두 금리는 잘 추적됩니다. 이는 ON RRP가 연방준비제도가 연방기금금리를 인상하는 데 도움이 될 수 있음을 의미합니다. 기억하십시오, 이 시기는 시스템에 막대한 양의 은행 준비금이 있었던 때입니다. 이러한 조건에서는 전통적인 공개시장운영이 작동하지 않았을 것입니다.
  \end{tcolorbox}
  \caption{ON RRP 금리와 연방기금금리 추이 설명}
  \label{fig:on_rrp_fed_funds_rate}
\end{figure}

\begin{summarybox}
  \textbf{핵심 정리: 풍부한 준비금 환경에서의 통화정책}
  \begin{itemize}
    \item 준비금이 많을 때, 전통적인 통화정책 도구는 이자율에 영향을 미치지 못할 수 있습니다.
    \item 연방준비제도는 연방기금금리에 하한선을 제공하기 위해 IORB와 ON RRP라는 새로운 도구를 만들었습니다.
  \end{itemize}
\end{summarybox}

\newpage
\section{빠르게 훑어보기 (1페이지 요약)}

\begin{tcolorbox}[summarybox, title=\textbf{중앙은행 통화정책 도구 요약}]
  \textbf{1. 공개시장운영 (OMO)}
  \begin{itemize}
    \item \textbf{개념:} 중앙은행이 증권(주로 국채)을 사고팔아 은행 준비금의 양을 조절.
    \item \textbf{목표:} 연방기금금리(Fed Funds Rate)를 목표 수준으로 유도.
    \textbf{작동:} 증권 매입 $\rightarrow$ 준비금 증가 $\rightarrow$ 연방기금금리 하락 (공급 곡선 우측 이동).
    \item \textbf{변화:} 2008년 금융위기 이후 준비금 풍부로 전통적 OMO 효과 감소.
  \end{itemize}

  \textbf{2. 재할인창구 (Discount Window)}
  \begin{itemize}
    \item \textbf{개념:} 중앙은행이 은행에 단기 현금 대출 제공 (최후의 대출자 역할).
    \item \textbf{목표:} 은행 유동성 지원, 뱅크런 방지, 연방기금금리 상한선 설정.
    \textbf{작동:} 재할인율(Discount Rate)을 연방기금금리 목표 상단에 설정. 연방기금금리가 재할인율 초과 시 은행은 재할인창구 이용 $\rightarrow$ 연방기금금리 하향 압력.
    \item \textbf{종류:} 주요 신용(Primary Credit)이 가장 중요 (건전한 은행 대상, 단기, 담보 대출).
  \end{itemize}

  \textbf{3. 비전통적 통화정책 (Unconventional Monetary Policy)}
  \begin{itemize}
    \item \textbf{등장 배경:} 2008년 금융위기 이후 단기 금리가 0에 가까워지면서 전통적 정책의 한계 발생.
    \item \textbf{포워드 가이던스 (Forward Guidance):}
      \begin{itemize}
        \item \textbf{개념:} 중앙은행이 미래 통화정책 의도(금리, 경제 전망 등)를 시장에 소통.
        \item \textbf{목표:} 시장 기대 형성, 불확실성 감소, 장기 금리 영향.
        \textbf{예시:} "상당히 오랜 기간" (extended period), 특정 경제 조건(실업률, 인플레이션)에 따른 금리 유지 약속 (조건부 가이던스), 닷 플롯(dot plots)을 통한 금리 전망 공개.
      \end{itemize}
    \item \textbf{대규모 자산 매입 (LSAPs / QE):}
      \begin{itemize}
        \item \textbf{개념:} 중앙은행이 장기 채권(국채, MBS 등)을 대량으로 매입.
        \item \textbf{목표:} 장기 채권 가격 상승 $\rightarrow$ 장기 수익률 하락 $\rightarrow$ 가계/기업 차입 비용 감소 $\rightarrow$ 투자/경제 활동 자극.
        \textbf{예시:} QE1, QE2, QE3, 오퍼레이션 트위스트, COVID-19 팬데믹 대응. 연준 대차대조표의 극적인 확장 초래.
      \end{itemize}
  \end{itemize}

  \textbf{4. 풍부한 준비금 환경에서의 통화정책}
  \begin{itemize}
    \item \textbf{문제점:} QE로 인해 은행 준비금이 과도하게 많아지면서, 전통적 OMO로는 금리 조절이 어려워짐.
    \item \textbf{지급준비금 이자 (IORB):}
      \begin{itemize}
        \item \textbf{개념:} 중앙은행이 은행이 보유한 준비금에 대해 이자를 지급.
        \item \textbf{목표:} 연방기금금리의 하한선 설정 (은행은 IORB보다 낮은 금리로 대출하지 않음).
      \end{itemize}
    \item \textbf{역환매조건부채권 시설 (ON RRP):}
      \begin{itemize}
        \item \textbf{개념:} 중앙은행이 비은행 금융기관(머니마켓 펀드 등)으로부터 증권을 담보로 돈을 빌리고 이자를 지급.
        \item \textbf{목표:} 비은행 기관의 단기 금리 하한선 설정 (ON RRP 금리보다 낮은 금리로 대출하지 않음).
      \end{itemize}
    \item \textbf{결과:} IORB와 ON RRP는 풍부한 준비금 환경에서도 연방기금금리를 효과적으로 통제할 수 있는 새로운 도구로 기능.
  \end{itemize}
\end{tcolorbox}

\end{document}
